\chapter{Introduction} \label{chapter1:introduction}
% \epigraph{It's really easy to make money if the tree grows long and skinny: compression is easier when betting on a skinny tree.}{Jack Lutz}

\section{Introduction} \label{section:Introduction}

These notes concern algorithmic information theory (AIT) and its application to geometric measure theory (GMT). This is motivated by the mounting evidence that mathematics under algorithmic restrictions offers insights into--and refinements to--the study of the fractal properties of sets. 

The origins of algorithmic information theory began in the 1960s with the formalization of \textit{Solomonoff induction}: a rule for selecting an hypothesis for the source from which some given data was produced. The rule substitutes a non-computable, universal prior distribution into Bayes' rule, and rewards algorithmically-simpler explanations \cite{Sol60}. R.~Solomonoff's conception of information content is most closely related to the algorithmic probability of producing a piece of data with respect to a universal Turing machine. 

An independent origin came just a few years later when A.~Kolmogorov defined the algorithmic (plain) complexity of individual objects, thereby refining the concept of classical Shannon information which only describes the mean information content across a distribution \cite{Kol68, Kol68b}. So, Kolmogorov's concept of information content was based on the degree of incompressibility with respect to a universal Turing machine. 

And shortly after, G.~Chaitin also independently introduced plain complexity and established its symmetry of information property, as well as establishing some initial notions of algorithmic randomness for infinite binary strings based on attaining maximal complexity along their prefixes \cite{Cha66}. 

P.~Martin-L{\"o}f characterized algorithmic randomness for an infinite binary sequence instead as the inability to cover it to arbitrary precision by effectively open statistical tests \cite{Mar66}. This adjusted A.~Church's original attempt: computable randomness, meant to achieve the intuitive property due to von~Mises that random sequences ought to pass all standard statistical tests such as the Law of Large Numbers. 

C.~Schnorr showed that both Chaitin's and Martin-L{\"o}f's notions of algorithmic randomness for infinite binary sequences were equivalent to yet another: those sequences on which any effective betting strategy (i.e., \textit{constructive martingale}) would fail to generate arbitrary gains when betting on subsequent bits \cite{Sch71}. Thus, what is now commonly known as \textit{Martin-L{\"o}f randomness} is considered robust for its equivalent characterizations in terms of Chaitin's notion of incompressibility, Martin-L{\"o}f's notion of typicality, and Schnorr's notion of unpredictability.

Moreover, L.~Levin confirmed the equivalence between Solomonoff, Kolmogorov, and Chaitin's approaches to quantifying algorithmic information content \cite{Lev73}. Levin's Coding Theorem equated Solomonoff's universal prior with $2^{-K(\cdot)}$, where $K(\cdot)$ denotes \textit{prefix-free} or \textit{self-delimiting Kolmogorov complexity}, or the minimal description-length of some data with respect to a universal, prefix-free Turing machine.

Whether viewed as the branch of theoretical computer science dedicated to the properties of computably-generated objects, or as a refinement of classical Shannon information theory by measuring the information content of individual objects, or as a project to formalize randomness: algorithmic information theory clearly offers some extra structure to otherwise ``classical'' mathematics, or the mathematics unconstrained by effectivity. 

Just as any natural, geometric object has an assignable geometric dimension that corresponds to the right ``exponential factor'' by which to measure its volume, the goal of many fractal dimension notions is to assign to a set an exponent appropriate to its geometry \cite{DMR05}. For instance, the \textit{Hausdorff dimension} of a set is roughly defined as the critical value of $s$ for which the volume of the set as measured using exponent $s$ switches from being infinite to null. Thus, Hausdorff dimension is naturally associated to the family of \textit{$s$-dimensional Hausdorff outer measures $\mathcal{H}^s$} (defined in Section~\ref{section:Measures_and_Hausdorff_Dimension}). These outer measures refine Lebesgue measure on null sets based on their covering properties. Yet, a further refinement to $\mathcal{H}^s$-nullness is exhibited by \textit{partial randomness}.

Generalizing the constructive martingales of Schnorr, J.~Lutz introduced \textit{constructive $s$-supergales}, where $s \geq 0$ is a dimension parameter \cite{Lut00}. Whenever a constructive $s$-supergale could make arbitrary gains betting on the subsequent bits of a given sequence (i.e., \textit{succeed} on the sequence), then that sequence would not be considered random with respect to $s$. Then, Lutz defined the \textit{constructive dimension} of a sequence to roughly be the minimal $s$ on which some constructive $s$-supergale could succeed. 

Shortly after, E.~Mayordomo elucidated the relationship between Lutz's constructive dimension and Kolmogorov complexity. In particular, the constructive dimension of a sequence matches L.~Staiger's limit inferior of the normalized Kolmogorov complexities of the sequence's prefixes \cite{May02, Sta89}. 

Both Lutz's and Mayordomo's contributions were built on decades of prior work toward understanding Kolmogorov complexity, including bounds on plain complexity by Shannon entropy \cite{Sta81}, as well as work by both B.~Ryabko \cite{Rya93, Rya94} and Staiger \cite{Sta98} partially relating the limit inferior of the normalized Kolmogorov complexity of prefixes to the speed at which a computable martingale could make arbitrary gains on a given sequence. 

Generalizing the work of Schnorr and Chaitin, K.~Tadaki introduced notions of Martin-L{\"o}f tests and partial randomness tailored to a dimension parameter $s \geq 0$ \cite{Tad02}, which in turn provided yet another characterization of constructive dimension as roughly the minimal value of $s$ for which a set could be covered by one of these tests. Thus, this effective notion of dimension serves as a refinement to Martin-L{\"o}f randomness, and similarly may be characterized in terms of incompressibility, typicality, and unpredictability. 

In effect, effective dimension quantifies just \textit{how random} a sequence behaves in the limit. In particular, it is possible for singleton sets to have nonzero effective dimension (in fact, most have full effective dimension). These would correspond to points which are not covered by a Martin-L{\"o}f-test appropriate for some dimension $s > 0$; or, equivalently, those admitting highly incompressible prefixes; or, those on which any constructive $s$-supergale would not succeed at making arbitrary gains for some $s > 0$.

In the recent past, computability theorists and geometric measure theorists alike have used some notions of effective dimension to prove classical results about the fractal dimension of certain sets. The so-called \textit{Point-to-Set Principles} by J.~Lutz and N.~Lutz have almost always served as the bridge between effective and ``classical'' results about Hausdorff or packing dimensions in Euclidean space \cite{LL18}. For example, their formula in Theorem~\ref{theorem:PTS_for_H_dimension} characterizes the Hausdorff dimension of a set in a pointwise manner: a strange result given that all singleton sets are classically zero-dimensional. But, as we have discussed, most points are not forthright in revealing their own smallness by effective means. Essentially, sets of large fractal dimension must contain many sufficiently random points.

There are still many open questions in regards to the potential applications of algorithmic information theory to geometric measure theory, as well as to what settings AIT applies. 

Note that this introduction as well as most of the contents feeding these notes were originally produced in my doctoral dissertation from 2025. 

\begin{itemize}
    \item \textbf{Chapter~\ref{chapter1:introduction}}: We fix some notation and review the basic definitions and results relevant to our discussion about AIT and GMT. Central to our work are the outer measures induced by premeasures on a space, various cover and test notions from algorithmic randomness, and measures of algorithmic information content such as Kolmogorov complexity and semimeasures. In particular, Theorem~\ref{theorem:continuous_semimeasure_is_supergale_is_strong_Solovay_tests} relates three notions from the partial randomness literature: success by a lower-semicomputable continuous semimeasure, success by a constructive supergale, and covering by a sequence of strong Solovay tests with rapidly shrinking weights.
    \item \textbf{Chapter~\ref{chapter:robustness_conditional_complexity}}: We establish some basic properties of the lift of Kolmogorov complexity to Euclidean space. Theorem~\ref{theorem:maxmin_and_minmax_conditional_complexity_equal} shows that the order of optimizers in the definition of conditional prefix complexity makes no difference up to sub-linear factors. Theorem~\ref{theorem:SOI_conditional_complexity_r_s_t} proves a conditional symmetry of information for conditional prefix complexity. Lemma~\ref{lemma:forward_data_processing_conditional} confirms a modulus processing inequality for conditional prefix complexity under computable, uniformly continuous maps, which implies conditional dimension is invariant under bi-computable, bi-Lipschitz continuous maps in Theorem~\ref{theorem:conditional_under_biLipschitz}.
\end{itemize}


\section{Notation} \label{section:Notation}

Let $\mathbb{R}$ denote the set of real numbers, $\mathbb{Q}$ the set of rational numbers, and $\omega$ the set of natural numbers (including $0$). Let $\mathbb{D} \subset \mathbb{Q}$ denote the set of all dyadic rational numbers (i.e., numbers of the form $k / 2^l$ for integers $k, l \geq 0$), and $\mathbb{D}_r$ the set of all $r$-dyadic rational numbers (i.e., all simplified dyadic rationals where the power of $2$ in the denominator is no more than $r$). 
% Call $\frac{1}{r} \mathbb{Z}$ the set of all rationals $q$ such that $r \cdot q \in \mathbb{Z}$. 
Let $2^n$ denote the set of length-$n$ binary strings, $2^{\leq n}$ the set of all binary strings length no greater than $n$, $2^{< \omega}$ the set of finite binary strings of any length, and analogous definitions for each of $2^{\geq n}$, $2^{> n}$, $2^\omega$, and $2^{\leq \omega}$. Similarly, $\omega^{<\omega}$ will denote the set of finite sequences of natural numbers, and $\omega^\omega$ the set of infinite sequences of natural numbers. Generally, for any space $X$, interpret $X^m$ as the set of length-$m$ tuples over $X$. In particular, elements of $(2^{\leq n})^m$ are $m$-tuples of binary strings of a common length less than or equal to $n$. 

If $\delta > 0$ and $A$ is a subset of a metric space $(\Omega, d)$, then let us define the \textit{$\delta$-neighborhood of $A$} to be $B_\delta(A) := \set{x \in \Omega}{d(x, A) < \delta}$. Then for any $x \in \Omega$, the set $B_\delta(x) := B_\delta(\seq{x})$ is the open ball of radius $\delta$ about $x$.

The cardinality of a set $A$ is denoted by $\abs{A}$. We also use $\abs{\cdot}$ for the absolute value on $\mathbb{R}$. Use $\lambda$ to denote the Lebesgue measure on both $2^\omega$ and $\mathbb{R}$. For any $p \in [0, \infty]$, $\norm[p]{\cdot}$ denotes the $L^p$-norm on Euclidean space, and we may simply write $\norm{\cdot}$ for the standard $L^2$-norm unless specified otherwise. While $\log_b$ will denote the base-$b$ logarithm for any $0 < b \neq 1$, we will always take $\log$ to mean the base-$2$ logarithm. Also let $\mathcal{Q}^m$ be the collection of all the $m$-dimensional dyadic cubes in $\mathbb{R}^m$, consisting of the sets $Q(z, a, r)$ of the form:
\begin{align*}
    [z_0 + a_0 \cdot 2^{-r}, z_0 + (a_0 + 1) \cdot 2^{-r}) \times \cdots \times [z_{m - 1} + a_{m - 1} \cdot 2^{-r}, z_{m - 1} + (a_{m - 1} + 1) \cdot 2^{-r}),
\end{align*}
where $z \in \mathbb{Z}^m$, $a \in \omega^m$, and $r \in \omega$. Collect into $\mathcal{Q}^m_r$ all the $m$-dimensional dyadic cubes $Q(z, a, r)$ with side-length $2^{-r}$.

It will often happen that certain equalities hold only up to an additive or multiplicative constant. For this reason we use the following notation: a statement of the form $f(x) \leq^+ g(x)$ for all $x \in X$ means that there exists a constant $c \in \mathbb{R}$ such that $f(x) \leq g(x) + c$ for all $x \in X$. Similarly $f(x) =^+ g(x)$ for all $x \in X$ means that $f(x) \leq^+ g(x) \leq^+ f(x)$ for all $x \in X$. Similarly, $f(x) \leq^* g(x)$ for all $x \in X$ will mean there exists a constant $c \in \mathbb{R}$ such that $f(x) \leq c \cdot g(x)$ for all $x \in X$, and similar for $f(x) =^* g(x)$. Similarly, the notation $f = O_{\boldsymbol{v}}(g)$ is to be interpreted as $f \leq^* g$ everywhere, where the implicit multiplicative constant may depend on the components in $\boldsymbol{v}$. And the notation $o(r)$ is used to represent a \textit{sub-linear term}, or a term which satisfies $\lim_{r \to \infty} \frac{o(r)}{r} = 0$.

The length of a string $\sigma \in 2^{\leq \omega}$ may be denoted by $\operatorname{len}(\sigma)$ (infinite whenever $\sigma \in 2^\omega$). A string $\tau \in 2^{\leq \omega}$ \textit{extends} another string $\sigma \in 2^{\leq \omega}$ if $\tau(n) = \sigma(n)$ for all $n < \operatorname{len}(\sigma)$, and this is denoted by $\sigma \preceq \tau$. Moreover, $\tau$ is a \textit{proper extension} of $\sigma$ if $\sigma \preceq \tau$ yet $\sigma \neq \tau$. Two strings $\sigma$ and $\tau$ are \textit{comparable} (denoted $\sigma \mid\mid \tau$) if either $\sigma \preceq \tau$ or $\tau \preceq \sigma$. String concatenation is denoted by $\sigma^\smallfrown \tau$ (or sometimes $\sigma \tau$, omitting the concatenation symbol) for any strings $\sigma \in 2^{<\omega}$ and $\tau \in 2^{\leq \omega}$. Formally, this is defined as the unique string of length $\operatorname{len}(\sigma) + \operatorname{len}(\tau)$ such that $\sigma \preceq \sigma^\smallfrown \tau$ and $(\sigma^\smallfrown \tau)(\operatorname{len}(\sigma) + n) = \tau(n)$ for all $0 \leq n < \operatorname{len}(\tau)$. 

A collection of finite binary strings $S \subseteq 2^{<\omega}$ is called \textit{prefix-free} if no two elements of $S$ are comparable, i.e., $(\forall \sigma, \tau \in S)[\neg(\sigma \mid\mid \tau)]$. Then for any $S \subseteq 2^{<\omega}$, denote by $\widehat{S}$ the collection of all minimal elements of $S$ with respect to the prefix partial relation, i.e.,
\begin{align*}
    \widehat{S} := \set{\sigma \in S}{\neg (\exists \tau \preceq \sigma)[\tau \in S]}.
\end{align*}
Notice that $\widehat{S}$ will always be a prefix-free subset of $S$. 

For any finite binary string $\sigma \in 2^{<\omega}$, denote by $[\sigma]$ the \textit{cylinder set above $\sigma$} comprising all infinite extensions of $\sigma$. Collect into $\mathcal{B}$ all of the cylinder sets above the finite strings from $2^{<\omega}$. Then $\mathcal{B}$ serves as a clopen basis for the product topology on $2^\omega$, in which every open set is of the form $[S] := \bigcup_{\sigma \in S} [\sigma]$, where $S \subseteq 2^{<\omega}$. Also, let $\bbrk{\sigma} := \set{\tau \in 2^{<\omega}}{\sigma \preceq \tau}$ denote the collection of finite binary strings extending $\sigma$. The truncation of a binary string $\sigma \in 2^{\leq \omega}$ up to the first $r \geq 0$ places is denoted by $\restrict{\sigma}{r} = (\sigma(i) : i < \min\seq{r, \operatorname{len}(\sigma)})$. 

The length of a string tuple $\sigma = (\sigma_1, \sigma_2, ..., \sigma_m)$ is defined as 
\begin{align*}
    \operatorname{len}(\sigma) := \min \seq{\operatorname{len}(\sigma_1), \operatorname{len}(\sigma_2), ..., \operatorname{len}(\sigma_m)},
\end{align*}
or infinite when all components are of infinite length. We may extend the prefix relation to string tuples $\sigma, \tau \in (2^{\leq \omega})^m$ component-wise, 
\begin{align*}
    \sigma \preceq \tau :\iff (\sigma_1 \preceq \tau_1) \AND (\sigma_2 \preceq \tau_2) \AND \cdots \AND (\sigma_m \preceq \tau_m).
\end{align*}
The truncation of a real number $x \in \mathbb{R}$ up to precision-level $r \in \omega$ may be defined as $\restrict{x}{r} := 2^{-r} \cdot \floor{x \cdot 2^r}$. Under the standard association between real numbers and their infinite binary expansions (under the convention of only permitting infinite tails of zeros), these two notions of truncation agree in the following sense. Without loss of generality, consider $x \in [0, 1)$ and let $x$ have the infinite binary expansion $0.x(0)x(1)x(2) \cdots$. Then, for any $r \in \omega$, it holds that its length-$r$ truncation as a string (followed by an infinite tail of zeros) $0.x(0) x(1) \cdots x(r - 1)^\smallfrown 0^\omega$ is the infinite binary expansion of the real number $\restrict{x}{r}$. 

Truncation for a real tuple $x = (x_1, x_2, ..., x_m) \in \mathbb{R}^m$ is taken component-wise:
\begin{align*}
    \restrict{x}{r} := \p{\restrict{x_1}{r}, \restrict{x_2}{r}, ..., \restrict{x_m}{r}} \in B_{2^{-r} \sqrt{m}}(x) \cap \mathbb{D}_r^m.
\end{align*}

Consider the collections of natural numbers: $\omega$, finite binary strings: $2^{<\omega}$, or dyadic rationals: $\mathbb{D}$. For the purposes of computability theory, we assume that between each pair of these sets there exists a computable bijection. Other collections can be added to this list, including the cylinder sets: $\mathcal{B}$, all the finite subsets of $\omega$, the integers: $\mathbb{Z}$, the rationals: $\mathbb{Q}$, the collection of finite sequences of natural numbers: $\omega^{<\omega}$, and finite-products of such sets. We say that each such collection is composed of \textit{finitary} objects, also known as \textit{hereditarily finite objects}. Whenever a definition is appropriate for any finitary object, we will state the definition using the symbol $\mathscr{N}$ as a placeholder for any one of these collections. 

For instance, given a standard collection of finitary objects $\mathscr{N}$ as described above and an element $n \in \mathscr{N}$, one might evaluate $\operatorname{len}(\cdot)$, truncation, or the prefix relation on the finite binary string to which $n$ is associated under a canonical, computable bijection between $\mathscr{N}$ and $2^{<\omega}$. 

Finitary objects often give rise to \textit{infinitary} objects via some limiting process and thus admit a notion of extension. For instance, any infinite binary string $x \in 2^\omega$ is the limiting sequence of its prefixes $\restrict{x}{r} \in 2^{<\omega}$ for $r \in \omega$. Any real $x \in \mathbb{R}$ is definable by Dedekind cuts over the rationals $\mathbb{Q}$. Any infinite subset of $\omega$ is the countable union of finite, nested subsets of $\omega$. Mappings between spaces of infinitary objects depend on the context, and each space of infinitary objects can take on extra metric or topological structure, so we will treat these spaces individually. Other examples include the set of countably-infinite sequences of natural numbers: $\omega^\omega$, and the set of infinite subsets of $\omega$ and their characteristic functions.

\section{Measures and Hausdorff Dimension} \label{section:Measures_and_Hausdorff_Dimension}

\subsection{Measures}

We begin by reviewing premeasures and their associated outer measures. For more details, refer to \cite{Rei08} and \cite{Rog70}.

\begin{definition}
    Let $\mathcal{C}$ be a collection of subsets of a set $\Omega$ containing $\emptyset$. Then a \textit{premeasure} on $\mathcal{C}$ is a function $\rho : \mathcal{C} \to [0, +\infty]$ and satisfies $\rho(\emptyset) = 0$.
\end{definition}

In the setting of Cantor space, it is customary to consider a premeasure on $2^\omega$ to be any non-negative mapping defined on the cylinder sets $\rho : \mathcal{B} \to [0, +\infty]$. Of course, such maps are in one-to-one correspondence with maps defined on the finite binary strings $\rho: 2^{<\omega} \to [0, +\infty]$. We may freely interchange between these two types by identifying any $\sigma \in 2^{<\omega}$ with its cylinder set $[\sigma] \in \mathcal{B}$.

% Introduce Method II from Rogers for premeasure -> outer measure
In metric spaces, there is a standard, general method (introduced as ``Method II'' in \cite{Rog70}) to induce an outer measure from a given premeasure. 

\begin{theorem}[Theorem 15 of \cite{Rog70}]
    Let $\rho$ be a premeasure on a collection $\mathcal{C}$ of subsets of a metric space $(\Omega, d)$. Define for each $X \subseteq \Omega$ and $\delta > 0$:
    \begin{align*}
        \mathcal{H}^\rho_\delta(X) &:= \inf_{(C_i)_{i \in \omega} \subseteq \mathcal{C}} \set{\sum_{i} \rho(C_i)}{\diam_d(C_i) \leq \delta, \bigcup_i C_i \supseteq X}, \quad \text{and} \quad \mathcal{H}^\rho(X) := \sup_{\delta > 0} \mathcal{H}^\rho_\delta(X).
    \end{align*}
    Then, $\mathcal{H}^\rho$ is an outer measure on $\Omega$.
\end{theorem}

We call $\mathcal{H}^\rho$ the \textit{Method II outer measure associated to the premeasure $\rho$}. And for each $\delta > 0$, we also call $\mathcal{H}^\rho_\delta$ the \textit{$\rho$-dimensional $\delta$-content} measure.

In \cite{Rog70}, C.~Rogers also defines an alternative ``Method I'' which differs from Method II in that it does not enforce any upper bound on the diameter of the cover elements used in the infimum:
\begin{align*}
    \mathcal{H}^\rho_{\operatorname{I}}(X) &:= \inf_{(C_i)_{i \in \omega} \subseteq \mathcal{C}} \set{\sum_{i} \rho(C_i)}{\bigcup_i C_i \supseteq X}.
\end{align*}
This outer measure is also sometimes denoted $\mathcal{H}^\rho_\infty$. While this method also produces an outer measure from a premeasure $\rho$ which agrees on null sets with $\mathcal{H}^\rho$, Rogers points out that only the results of Method II can be guaranteed to be additive on sets which are separated by positive distance \cite{Rog70}. We proceed with Method II as our preferred way to produce outer measures from premeasures.

Cantor space: $2^\omega$, is a metric space when imbued with the standard metric,
\begin{align*}
    d(x, y) := \begin{cases}
        2^{-N} & \text{if $x \neq y$ and $N = \min\set{n \in \omega}{x(n) \neq y(n)}$},\\
        0 & \text{if $x = y$}.
    \end{cases}
\end{align*}
Note that this metric is compatible with the product topology generated by the cylinder sets in $\mathcal{B}$. 

\subsection{Hausdorff Measures}

An important class of outer measures on any metric space $(\Omega, d)$ are the \textit{Hausdorff outer measures}, which are those Method II outer measures induced from premeasures whose values depend only on the diameter of the set. First, declare a \textit{dimension function} to be any non-negative, non-decreasing, continuous-on-the-right function $h$ defined on all non-negative real numbers and satisfying $h(t) = 0 \iff t = 0$. One may associate to any dimension function $h$ its \textit{Hausdorff premeasure} defined on all $X \subseteq \Omega$ as follows:
\begin{align*}
    \rho_h(X) := (h \circ \diam_d)(X).
\end{align*}
In Cantor space, $\rho_h(\sigma) = h\p{2^{- \operatorname{len}(\sigma)}}$ for any $\sigma \in 2^{<\omega}$. So, any Hausdorff premeasure on $2^\omega$ is \textit{length-invariant}, in that $\rho_h(\sigma)$ only depends on the length of $\sigma$. Later, we will briefly focus on \textit{convex} Hausdorff premeasures $\rho$ over Cantor space, which satisfy:
\begin{align*}
    \rho(\sigma) \leq \rho(\sigma^\smallfrown 0) + \rho(\sigma^\smallfrown 1),
\end{align*}
for all $\sigma \in 2^{<\omega}$. It is easy to see that all length-invariant, convex Hausdorff premeasures over $2^\omega$ are also \textit{strongly convex}, meaning for all $\sigma \in 2^{<\omega}$ and $i \in \seq{0, 1}$, 
\begin{align*}
    \rho(\sigma ^\smallfrown i) \geq \frac{\rho(\sigma)}{2}.
\end{align*}

Essential to fractal geometry is the family of $s$-dimensional Hausdorff premeasures $\rho_s := \rho_{h_s}$, where $h_s : t \mapsto t^s$ is the dimension function associated to $s \geq 0$. If $\delta > 0$, let $\mathcal{H}^s_{\delta} := \mathcal{H}^{\rho_s}_\delta$ denote the \textit{$s$-dimensional Hausdorff $\delta$-content}, and $\mathcal{H}^s := \mathcal{H}^{\rho_s}$ the \textit{$s$-dimensional Hausdorff outer measure}. 

For any outer measure $\mu$ on a metric space $(\Omega, d)$, we say a set $X \subseteq \Omega$ is \textit{$\mu$-null} if $\mu(X) = 0$. Any $\mathcal{H}^s$-null $X$ is also called \textit{$s$-null}. 

Suppose $\mathcal{H}^s(X) < \infty$ for some $s \geq 0$. This implies that $X$ is $s_+$-null for any $s_+ > s$ too. Alternatively, if $\mathcal{H}^s(X) > 0$ for some $s \geq 0$, then it must be that $\mathcal{H}^{s_-}(X) = \infty$ for all $0 \leq s_- < s$. So, the shape of the function $s \mapsto \mathcal{H}^s(X)$ is either constant with value $0$ or $\infty$, or a step from $\infty$ to $0$ (where, at the critical value of $s$ where the step occurs, $\mathcal{H}^s(X)$ could take on any value in $[0, \infty]$). This motivates the following definition for the \textit{Hausdorff dimension} of a set $X$ \cite{Hau19, Fal90}:
\begin{align}
    \label{equation:definition_of_Hausdorff_dimension}
    \dim_{\operatorname{H}}(X) := \inf\set{s \geq 0}{\text{$X$ is $s$-null}},
\end{align} 
or $+\infty$ if the infimum is taken over the empty set.

\section{Computability Theory} \label{section:Computability_Theory}

\subsection{Computability by Machines}

% Church-Turing thesis

We briefly review the necessary components of computability theory for discussing algorithmic randomness. In this dissertation, we accept Turing's model of computation as the ``right'' way to characterize \textit{algorithmic} or \textit{effective} procedures. According to the Church-Turing thesis, the subsets of the natural numbers or functions from natural numbers to natural numbers which are ``intuitively algorithmic'' are exactly those which may be computed by Turing machines. See \cite{DH10} for more on the formalism of Turing machines. 

To begin, we will consider partial functions of the form $\Phi : \subseteq \omega \to \omega$. If $\operatorname{dom}(\Phi) = \omega$, call $\Phi$ \textit{total}. Each Turing machine $M$ \textit{computes} a partial function $\Phi_M : \subseteq \omega \to \omega$ satisfying $n \in \operatorname{dom}(\Phi_M) \iff M(n) \downarrow$, and $n \in \operatorname{dom}(\Phi_M)$ implies $M(n) \downarrow= \Phi_M(n)$. All such $\Phi_M$ are called \textit{partial computable (p.c.)} functions. Moreover, whenever $A = \operatorname{dom}\Phi_M$ is the domain of a Turing machine $M$, then $M$ is said to \textit{computably enumerate} $A$, making $A$ \textit{computably enumerable (c.e.)}. If both $A$ and its complement $\omega \setminus A$ are c.e., then $A$ is called \textit{computable}. 

Proposition 2.2.2 of \cite{DH10} helps to motivate the name ``computably enumerable.'' It states that $A$ is c.e.~if and only if either $A$ is empty or there is a total computable function $\Phi$ from $\omega$ onto $A$. That is, the non-empty c.e.~sets are exactly those enumerated in some order by a Turing machine.

By the Enumeration Theorem (see Theorem 2.1.2 of \cite{DH10}), there exists a single algorithm to enumerate all partial computable functions $(\Phi_e)_{e \in \omega}$. Thus, there exists a \textit{universal} partial computable function $f$ of two variables such that $f(e, n) = \Phi_e(n)$ for all $e, n \in \omega$ (note that they could both diverge). Let $\mathbf{U}$ be a Turing machine computing $f$, which is called a \textit{universal Turing machine}. 

% Oracle Turing machines, and existence of a universal oracle machine
A similar result holds for oracle machines or Turing functionals (see \cite{RS15} and \cite{DH10} for more details). Intuitively, oracle machines are Turing machines which receive and may make use of the extra information encoded in some \textit{oracle}. An oracle machine $M$ given oracle $B \in 2^{\leq \omega}$ is said to compute a \textit{$B$-partial computable function} $\Phi^B_M : \subseteq \omega \to \omega$ as before. Similarly, there is a single oracle machine (also called $\mathbf{U}$) which, given an oracle $B \in 2^{\leq \omega}$, will enumerate all $B$-partial computable functions $(\Phi^B_e)_{e \in \omega}$.

% Describe what "relativization" means in computations
If there is an oracle machine $M$ that computes the set $A \subseteq \omega$ given an oracle $B$, then $A$ is said to be \textit{Turing reducible} to $B$, or \textit{$B$-computable}. This is denoted by $A \leq_{\operatorname{T}} B$. Similarly, a set $A$ is \textit{computably enumerable in $B$} (or, \textit{$B$-c.e.}) whenever $A = \operatorname{dom} \Phi$ for some $B$-partial computable function $\Phi$. \textit{Relativization} is the process of generalizing notions or results from Turing machines to oracle machines given arbitrary oracle-power. Most results we cite from computability theory may be relativized appropriately. 

In later sections, we will also make use of lower- or upper-semicomputability. A real number $x \in \mathbb{R}$ is \textit{lower-semicomputable} (or \textit{left-c.e.}) if its Dedekind left cut $\set{q \in \mathbb{Q}}{q < x}$ is c.e. Similarly, a real-valued map $f: \omega \to \mathbb{R}$ is \textit{lower-semicomputable} if its lower graph $\set{(n, q) \in \omega \times \mathbb{Q}}{q < f(n)}$ is c.e. If $f : \omega \to [a, \infty)$ has a known, computable codomain which is bounded from below, we may equivalently characterize lower-semicomputability by the ability to be approximated from below by some computable \textit{left-approximator}. That is, suppose there exists a computable map $\widehat{f} : \omega \times \omega \to [a, b]$ satisfying:
\begin{align*}
    \widehat{f}(i, r) \leq \widehat{f}(i, r + 1) \leq f(i), \quad \text{and} \quad \lim_{r \to \infty} \widehat{f}(i, r) = f(i),
\end{align*}
for all $i, r \in \omega$. Then, it is straightforward to argue that the lower graph of $f$ is computably enumerable. Conversely, from an enumeration of the lower graph of $f$, one may define such a left-approximator $\widehat{f}$ to $f$: for a given $i \in \omega$, one may approximate $f(i)$ as the largest rational $q$ such that $(i, q)$ has been enumerated into the lower graph of $f$. If no such rational has yet been approximated, we may take $\widehat{f}(i, r) = a$. Analogous definitions work for upper-semicomputability. 

% Insert here that these notions of computability may be generalized to other domains of finitary domains from the noatation section through a computable bijections between that collection and the natural numbers. 
The definitions above are given for either subsets of the natural numbers or functions on the natural numbers, but they may be extended to the other standard collections of finitary objects discussed in Section~\ref{section:Notation} via a computable bijection.

\subsection{Computable Analysis}

Now, let us import some notions from computable analysis as presented in \cite{KF82}. For the rest of this section, fix the ambient dimensions $m, n \in \omega$ and an oracle machine $M$. 

Fix a map $\phi : \omega \to \mathbb{D}^m$ from the natural numbers to dyadic rationals in $\mathbb{R}^m$ satisfying $\phi(r) \in \mathbb{D}^m_r$ for all $r \in \omega$. Then, $\phi$ is said to be a \textit{Cauchy representation} of $x \in \mathbb{R}^m$ if $\abs{\phi(r) - x} \leq 2^{-r}$ for all $r \in \omega$. If $\phi$ is a Cauchy representation of $x$ further satisfying $\phi(r) \leq x < \phi(r) + 2^{-r}$ for all $r \in \omega$, then $\phi$ is said to be a \textit{standard Cauchy representation} for $x$. Such a $\phi$ can compute the binary expansion of $x$ having no infinite tail of ones. For any $\phi: \omega \to \mathbb{D}^m$, we will denote by $M^\phi$ the partial map computed by $M$ given $\phi$ encoded as an infinite binary string for its oracle. And $M^x$ will denote the partial map computed by $M^\phi$, where $\phi$ is some standard Cauchy representation of $x$. 

$M$ is said to \textit{compute} a real function $f: \mathbb{R}^m \to \mathbb{R}^n$ if $\phi$ being a Cauchy representation for some $x \in \dom(f)$ implies $M^\phi$ is a Cauchy representation of $f(x)$.

Now, fix a metric space $(\Omega, d)$. Two subsets $X$ and $Y$ of $\Omega$ are said to be \textit{$\varepsilon$-close} for $\varepsilon > 0$ if for each $x \in X$, there exists $y \in Y$ such that $d(x, y) < \varepsilon$, and vice versa. The \textit{Hausdorff distance} between $X$ and $Y$ is then defined:
\begin{align*}
    d_{\operatorname{H}}(X, Y) = \inf \set{\varepsilon > 0}{\text{$X$ is $\varepsilon$-close to $Y$}}.
\end{align*}

We may now recall the effective versions of certain metric spaces as presented in \cite{IK21}. 

The triple $\mathbf{\Omega} = (\Omega, d, \alpha)$ is a \textit{computable metric space} if $\alpha = (\alpha_i)_{i \in \omega}$ is a dense sequence in $(\Omega, d)$ for which the function mapping $(i, j) \mapsto d(\alpha_i, \alpha_j)$ is computable. 

Take $\mathbf{\Omega} = (\Omega, d, \alpha)$ to be a computable metric space. If $d$ happens to be a complete metric, call $\mathbf{\Omega}$ a \textit{computable Polish space}. Any subset $X \subseteq \Omega$ of a computable metric space is called \textit{computably compact} if either $X = \emptyset$ or there exists a computable map $f : \omega \to \omega$ such that for each $k \in \omega$,
\begin{align*}
    d_{\operatorname{H}}(X, \Lambda_{f(k)}) < 2^{-k},
\end{align*}
where $\Lambda_e$ is the $e$-th collection enumerated in a standard, computable enumeration of all finite sub-collections of $\alpha$. Then, the computable metric space $\mathbf{\Omega}$ is called \textit{effectively compact} if $\Omega$ is computably compact in $\mathbf{\Omega}$.

\subsection{Lightface Hierarchies}

% Arithmetic hierarchy 
Kleene's \textit{arithmetic hierarchy} assigns definability-complexity levels to subsets of the natural numbers. Any set $A \subseteq \omega$ is a \textit{$\Sigma^0_n$-set} if there is a computable relation $R$ on $n + 1$ arguments such that 
\begin{align*}
    a \in A \iff \exists x_1 \forall x_2 \cdots Qx_n R(x_1, ..., x_n, a),
\end{align*}
where the quantifiers alternate and end with $Q$ being $\exists$ whenever $n$ is odd, and $\forall$ when $n$ is even. A \textit{$\Pi^0_n$-set} is the complement of a $\Sigma^0_n$ set. And a set which is both $\Sigma^0_n$ and $\Pi^0_n$ is called $\Delta^0_n$. For convenience, each of $\Delta^0_0$, $\Sigma^0_0$, and $\Pi^0_0$ are defined as containing exactly the computable sets. Moreover, a set is computably enumerable if and only if it is $\Sigma^0_1$, and computable if and only if it is $\Delta^0_1$. This hierarchy may be relativized to any oracle by permitting the relation be Turing-reducible to it. 

% Lightface Borel Hierarchy for classes & Connection to Borel hierarchy by Descriptive Set Theory
Suppose $S \subseteq 2^{<\omega}$ is a c.e.~set of finite binary strings. Then $[S] := \bigcup_{\sigma \in S}[\sigma]$ is called a \textit{$\Sigma^0_1$-class}, or an \textit{effectively open} subset of Cantor space. The complement of a $\Sigma^0_1$-class is called a \textit{$\Pi^0_1$-class}, or an \textit{effectively closed} subset. Starting from these base classes, an analogous hierarchy to Kleene's arithmetic hierarchy forms for subsets of Cantor space. There are two ways to build higher levels of this lightface hierarchy. The first is by computable relations: a subset $U \subseteq 2^\omega$ is a \textit{$\Sigma^0_n$-class} for $n \in \omega$ if there is a computable relation $R$ on $n$ arguments such that
\begin{align*}
    \alpha \in U \iff \exists x_1 \forall x_2 \cdots Q x_n\ R(x_1, x_2, ..., \restrict{\alpha}{x_n}),
\end{align*}
again by alternating quantifiers with an appropriate choice for the innermost one $Q$. The complement of a $\Sigma^0_n$-class is a \textit{$\Pi^0_n$-class}. But we may equivalently conceive of the $\Sigma^0_{n + 1}$-classes as uniformly-computable unions of $\Pi^0_n$-classes. This may be formalized using effective Borel codes. And a set which is both a $\Sigma^0_n$-class and a $\Pi^0_n$-class is also called a \textit{$\Delta^0_n$-class}. These notions may be relativized to any oracle $B \in 2^{\leq \omega}$, notated with $B$ in the superscript: $\Sigma^{0, B}_n$, etc. The Fundamental Theorem of effective descriptive set theory states that for any $n \in \omega$, it holds that $\boldsymbol{\Sigma}^0_n = \bigcup_{B \in 2^{\leq \omega}} \Sigma^{0, B}_n$, where $\boldsymbol{\Sigma}^0_n$ denotes the corresponding level of the (boldface) Borel hierarchy from classical descriptive set theory. This motivates naming the hierarchy of $\Sigma^0_n$ and $\Pi^0_n$-classes as the \textit{lightface Borel hierarchy}. 

\section{Algorithmic Randomness} \label{section:Algorithmic_Randomness}
We review the tests-based approach to algorithmic randomness, all with respect to a given premeasure.

\subsection{Weights}

% Direct, prefix, and covering weight notions (all wrt to a measure)

\begin{definition}
    Let $\mathcal{V} \subseteq 2^{<\omega}$ be a collection of finite binary strings and $\rho$ be a premeasure on $2^\omega$. Define for $\mathcal{V}$ its:
    \begin{itemize}
        \item \textit{Direct $\rho$-weight}: $\operatorname{DW}_\rho(\mathcal{V}) := \sum_{\sigma \in \mathcal{V}} \rho([\sigma])$,
        \item \textit{Prefix $\rho$-weight}: $\operatorname{PW}_\rho(\mathcal{V}) := \sup \set{\operatorname{DW}_\rho(\mathcal{P})}{\mathcal{P} \subseteq \mathcal{V} \text{ is prefix-free}}$,
        \item \textit{Vehement $\rho$-weight}: $\operatorname{VW}_\rho(\mathcal{V}) := \inf \set{\operatorname{DW}_\rho(\mathcal{W})}{\mathcal{W} \subseteq 2^{< \omega} \text{ and } [\mathcal{V}] \subseteq [\mathcal{W}]}$.
    \end{itemize}
\end{definition}

When working with the $s$-dimensional Hausdorff premeasure $\rho_s$ for some $s \geq 0$, we might simply write $\operatorname{DW}_s(\mathcal{V})$ for the \textit{direct $s$-weight} of $\mathcal{V}$, and similar for the other weight notions.

For any $\mathcal{V} \subseteq 2^{<\omega}$, we defined $\widehat{\mathcal{V}}$ in Section~\ref{section:Notation} to include only those strings in $\mathcal{V}$ which are minimal with respect to the prefix relation. Then $\widehat{\mathcal{V}}$ is prefix-free and satisfies $[\widehat{\mathcal{V}}] = [\mathcal{V}]$. Using this fact, it is straightforward to verify that the above weight notions are ordered as follows:
\begin{align*}
    \operatorname{VW}_\rho(\mathcal{V}) \leq \operatorname{PW}_\rho(\mathcal{V}) \leq  \operatorname{DW}_\rho(\mathcal{V}).
\end{align*}

These covering weights also relate to the $s$-dimensional Hausdorff outer measures. In particular, if $\mathcal{H}^s(X) < c$ for some dimension parameter $s \geq 0$, constant $c > 0$, and subset $X \subseteq 2^\omega$, then it follows that for every $r \in \omega$, there exists a collection of finite binary strings $S \subseteq 2^{\geq r}$ such that $X \subseteq [S]$ and $\operatorname{DW}_s(S) < c$.

\subsection{Tests}
% ML, strong ML, Solovay, strong Solovay test notions 
We introduce some test notions, or effective schemes for identifying null-sets with respect to some premeasure. For the rest of the section, we will only consider upper-semicomputable premeasures on $2^\omega$. This restriction is made to ensure the collection of all tests of some type is uniformly-c.e. By relativizing these tests to some oracle, one could expand the class of premeasures considered to those which are upper-semicomputable in $B$.

\begin{definition}
    Let $\rho$ be an upper-semicomputable premeasure on $2^\omega$. Suppose $\mathcal{U} = (U_n)_{n \in \omega}$ is a uniformly-c.e.~sequence of collections $U_n \subseteq 2^{<\omega}$ of finite binary strings. Then define $\mathcal{U}$ to be: 
    \begin{itemize}
        \item a \textit{Martin-L{\"o}f-$\rho$-test} if $\operatorname{DW}_\rho(U_n) \leq 2^{-n}$ for all $n \in \omega$,
        \item a \textit{strong Martin-L{\"o}f-$\rho$-test} if $\operatorname{PW}_\rho(U_n) \leq 2^{-n}$ for all $n \in \omega$, and
        \item a \textit{vehement Martin-L{\"o}f-$\rho$-test} if $\operatorname{VW}_\rho(U_n) \leq 2^{-n}$ for all $n \in \omega$.
    \end{itemize}
    And if $\mathcal{V} \subseteq 2^{<\omega}$ is a c.e.~collection of finite binary strings, define $\mathcal{V}$ to be:
    \begin{itemize}
        \item a \textit{Solovay-$\rho$-test} if $\operatorname{DW}_\rho(\mathcal{V}) < + \infty$, and
        \item a \textit{strong Solovay-$\rho$-test} if $\operatorname{PW}_\rho(\mathcal{V}) < + \infty$.
    \end{itemize}
\end{definition}

% Order the test notions
We use $\mathcal{T}$ as a placeholder for any one of the above test notions. Once again, for any Hausdorff premeasure $\rho_s$ with $s > 0$ being left-c.e., we simplify each of the above terms as \textit{$\mathcal{T}$-$s$-tests}. And when $s = 1$, we might simply refer to them as \textit{$\mathcal{T}$-tests}. Furthermore, any one of these test notions may be relativized to an oracle $B \in 2^{\leq \omega}$, regarding computable-enumerability as in $B$. We may refer to such tests as \textit{$\mathcal{T}$-$B$-$\rho$-tests}. 

A Martin-L{\"o}f-$\rho$-test (or ML-$\rho$-test) is the strictest notion of test we have defined, meaning from any Martin-L{\"o}f-$\rho$-test $\mathcal{U}$, one may produce a $\mathcal{T}$-$\rho$-test for any of the other test notions $\mathcal{T}$. In fact, such a $\mathcal{U}$ is already both a strong and vehement ML-$\rho$-test, and generates the collection $\mathcal{V} := \bigcup_n U_n$ which is a (strong) Solovay-$\rho$-test. Of course, any strong ML-$\rho$-test is also a vehement ML-$\rho$-test, and any Solovay-$\rho$-test is already a strong Solovay-$\rho$-test.
    
\subsection{Covers} 

Tests act as effective covers of subsets of Cantor space. 

% Passing respective tests
\begin{definition}
    \label{definition:test_coverings}
    Let $\rho$ be an upper-semicomputable premeasure on $2^\omega$ and $X \subseteq 2^\omega$ be a subset. 
    \begin{itemize}
        \item Let $\mathcal{U} = (U_n)_{n \in \omega}$ be a (possibly strong or vehement) Martin-L{\"of}-$\rho$-test. Then $\mathcal{U}$ is said to \textit{cover} $X$ whenever $X \subseteq \bigcap_{n} U_n$. Otherwise, $X$ is said to have \textit{passed} $\mathcal{U}$.
        \item Alternatively, let $\mathcal{V}$ be a (possibly strong) Solovay-$\rho$-test. Then $\mathcal{V}$ is said to \textit{cover} $X$ whenever, for each $x \in X$, there are infinitely many $r \in \omega$ for which $\restrict{x}{r} \in \mathcal{V}$. Otherwise, $X$ is said to have \textit{passed} $\mathcal{V}$.
    \end{itemize}
\end{definition}

For a fixed test notion $\mathcal{T}$ (e.g., Martin-L{\"o}f) and upper-semicomputable premeasure $\rho$ on $2^\omega$, an infinite binary string $x \in 2^\omega$ is said to be $\mathcal{T}$-$\rho$-\textit{random} if $\seq{x}$ passes all $\mathcal{T}$-$\rho$-tests. Otherwise, $\seq{x}$ is covered by some such test, exhibiting itself to be (effectively) \textit{$\mathcal{T}$-$\rho$-null}. 

Note that if $\rho$ is an upper-semicomputable, Hausdorff premeasure (i.e., the premeasure associated to an upper-semicomputable dimension function $h$ as described in Section~\ref{section:Measures_and_Hausdorff_Dimension}), being $\mathcal{T}$-$\rho$-null for any of the test notions $\mathcal{T}$ defined here implies being $\mathcal{H}^\rho$-null. Suppose, for instance, that $\mathcal{U}$ were a vehement ML-$\rho$-test and $\delta > 0$. Then, it is possible to show that there exists a sequence $(W_n)_{n \in \omega}$ where each $W_n \subseteq 2^{<\omega}$ consists of strings of diameter no greater than $\delta$ and is such that $[U_n] \subseteq [W_n]$ and $\operatorname{DW}_\rho(W_n) \leq 2 \cdot 2^{-n}$. If $\mathcal{V} = \bigcup_n W_n$, then 
\begin{align*}
    \sum_{\sigma \in \mathcal{V}} \mathcal{H}_\delta^\rho([\sigma]) = \sum_n \sum_{\sigma \in W_n} \operatorname{DW}_\rho([\sigma]) = \sum_n \operatorname{DW}_\rho(W_n) \leq \sum_n 2 \cdot 2^{-n} < \infty.
\end{align*}
Since $\delta > 0$ was arbitrary, $\mathcal{H}^\rho$ is an outer measure, and $X \subseteq \limsup_{\sigma \in \mathcal{V}} [\sigma]$, the (first) Borel-Cantelli lemma implies $\mathcal{H}^\rho(X) = 0$. 
Additionally, supposing $\mathcal{V}$ were a strong Solovay-$\rho$-test covering $X$ and $\delta > 0$, then we could restrict $\mathcal{V}$ to only those strings of diameter no more than $\delta$ while maintaining $\mathcal{V}$ being a strong Solovay-$\rho$-test covering $X$. Then, for each $\varepsilon > 0$, it is possible to show that there exists a prefix-free subset $\mathcal{P} \subseteq \mathcal{V}$ such that $[\mathcal{V}] \subseteq [\mathcal{P}]$ and $\operatorname{DW}_\rho(\mathcal{P}) < \varepsilon$. Then, by monotonicity, $\mathcal{H}^\rho_\delta(X) \leq \operatorname{DW}_\rho(\mathcal{P}) < \varepsilon$. And since both $\delta, \varepsilon > 0$ were arbitrary, we conclude $X$ is $\mathcal{H}^\rho$-null.

For any upper-semicomputable premeasure $\rho$ on $2^\omega$ and any of the above test notions $\mathcal{T}$, it follows from the property that $\mathcal{H}^\rho$ is an outer measure that the set of non-$\mathcal{T}$-$\rho$-randoms is $\mathcal{H}^\rho$-null.

% Order the randomness notions
Weaker randomness notions come from stricter test notions. So, the class of Martin-L{\"o}f-$\rho$-randoms contains all the other classes of $\mathcal{T}$-$\rho$-randoms discussed here. 

% Note about randomness with respect to an outer measure... not necessary to include
% It is possible to go further and develop a notion of Martin-L{\"o}f randomness with respect to an outer measure on $2^\omega$ \cite{Rei08}. One such way involves setting a scheme for representing the premeasures on $2^\omega$ by Cauchy-sequences of some family of basic, computable premeasures. In that context, randomness with respect to an outer measure is defined as randomness relative to some representation for that measure. This extended notion of randomness will not be necessary for our discussion.

\subsection{More Covering Notions}
% Look into Katz & Tao, or Orponen's work, to give some basic facts about $(C, \delta, s)$-sets

Let us start by reviewing some definitions originally appearing in \cite{KT01} (and re-appearing slightly differently in \cite{Orp21}). 

\begin{definition}
    \label{definition:strongly_covers}
    A countable collection $(X_i)_{i \in \omega}$ of sets \textit{strongly covers} another set $X$ if each element of $X$ is also an element of infinitely many $X_i$. 
\end{definition}
Note that in Definition~\ref{definition:test_coverings}, we consider $X$ to be covered by a Solovay-type test if that test strongly covers $X$ in the sense of Definition~\ref{definition:strongly_covers}.

\begin{definition}
    For any (bounded) set $X \subseteq \mathbb{R}^m$ and $\delta > 0$, define,
    \begin{align*}
        N(X, \delta) := \min \set{\abs{Y}}{\text{$Y$ is a set of $\delta$-balls covering $X$}}.
    \end{align*}
\end{definition}

The quantity $N(X, \delta)$ is exactly what appears in the definitions of the box-counting dimensions, as it captures the least number of $\delta$-balls required to cover $X$ (see \cite{Fal90}). This fractal dimension always bounds the Hausdorff dimension from above.

\begin{definition}
    \label{definition:C_delta_s_set}
    Fix $m \in \omega$, and let $s, \delta > 0$ and $C \geq 1$. A finite set $P \subseteq \mathbb{R}^m$ is called a \textit{$(C, \delta, s)$-set} if for any $x \in \mathbb{R}^m$ and $\delta \leq \varepsilon \leq 1$,
    \begin{align*}
        \abs{P \cap B_\varepsilon(x)} \leq C \cdot \p{\frac{\varepsilon}{\delta}}^s.
    \end{align*}
\end{definition}
Suppose that $X \subseteq \mathbb{R}^m$ has $\dim_{\operatorname{H}} X = s$. Then any $(C, \delta, s)$-set will have no more points in $X$ than a constant times what any $\delta$-net on $X$ would. That is, $(C, \delta, s)$-sets are spread out such that they never dedicate many points to sets of small Hausdorff dimension.

Note that it is not vital that we use $r$-balls in these definitions for $N(X, \delta)$ or $(C, \delta, s)$-sets. For instance, replacing these with dyadic cubes of comparable side-lengths would produce roughly equivalent notions. % Falconer has an argument about this...

Over Euclidean space, there is a relationship between Hausdorff dimension and $(C, \delta, s)$-sets. For instance, a set of Hausdorff dimension less than $s$ will have a $(C, \delta, s)$-set whose $\delta$-neighborhood strongly covers the set. The precise result originally appeared in \cite{KT01} by N.~Katz and T.~Tao, but was simplified and slightly generalized by T.~Orponen in \cite{Orp21}.

\begin{lemma}[Lemma 7.5 of \cite{KT01}; Lemma 2.1 of \cite{Orp21}]
    \label{lemma:Katz_Tao}
    Let $0 < s \leq m$ and let $X \subseteq \mathbb{R}^m$ be a subset with $\dim_{\operatorname{H}} X < s$. Then, there exists a constant $C \geq 1$ depending only on $m$, $s$, and $\dim_{\operatorname{H}} X$ such that: for every $k \in \omega$, there exists a $(Ck^2, 2^{-k}, s)$-set $P_k$ such that the sequence $\p{B_{C_m \cdot 2^{-k}}(P_k)}_{k \in \omega}$ strongly covers $X$, where $C_m \geq 1$ only depends on $m$.
\end{lemma}

The above lemma was leveraged by Orponen to generalize a weak form of the Marstrand-Mattila Projection Theorem for arbitrary subsets of Euclidean space \cite{Orp21} (see Theorem~\ref{theorem:Orponen_weak_Marstrand}). We will present an effective proof of the Katz-Tao lemma later in Section~\ref{section:Effective_versions_of_Combinatorial_Proofs}.

% Optimal covers
Also recall the related notion of an \textit{optimal cover} originally by J.~Miller \cite{Mil11}. 
\begin{definition}
    \label{definition:optimal_cover}
    Let $\rho$ be a premeasure on $2^\omega$ and $S \subseteq 2^{<\omega}$. Then a \textit{$\rho$-optimal cover} of $S$ is a set $S^* \subseteq 2^{<\omega}$ satisfying $[S] \subseteq [S^*]$ and $\operatorname{DW}_\rho(S^*) = \operatorname{VW}_\rho(S)$.
\end{definition}

Optimal covers witness the vehement $\rho$-weight of a set. In general, it is unclear whether optimal covers exist. However, a result by P.~Hudelson proves a necessary condition for the existence of optimal covers. 

\begin{lemma}[Theorem 3.4.22 of \cite{Hud13}]
    \label{lemma:exists_optimal_cover_for_convex}
    Let $\rho$ be a convex Hausdorff premeasure on $2^\omega$. Then, for all $S \subseteq 2^{<\omega}$, there exists a $\rho$-optimal cover of $S$. 
    
    In particular, if $\rho$ is also computable and $S$ is c.e., then one may find another c.e.~set $A$ (uniformly in $\rho$ and $S$) such that $\widehat{A}$ is a $\rho$-optimal cover of $S$.
\end{lemma}

Essential to the proof of Lemma~\ref{lemma:exists_optimal_cover_for_convex} is the fact that whenever $\rho$ is convex and $\sigma \in 2^{<\omega}$, then $\seq{\sigma}$ is the $\operatorname{DW}_\rho$-minimal cover of $[\sigma]$ among all other collections $S \subseteq \bbrk{\sigma}$ covering $[\sigma]$. 

Observe that an optimal cover will necessarily be prefix-free by minimality. When $\rho = \rho^s$, simplify the term to an \textit{$s$-optimal cover}. Optimal covers were originally employed by Miller to construct a $\Delta^0_2$-definable real with effective dimension equal to $1 / 2$ yet unable to compute any real of greater effective dimension \cite{Mil11}. We will make use of optimal covers for effectivizing Lemma~\ref{lemma:Katz_Tao} in Section~\ref{section:Effective_versions_of_Combinatorial_Proofs}.

\section{Prefix Complexity} \label{section:Prefix_Comlexity}

\subsection{Incompressibility}

Recall the universal oracle machine $\mathbf{U}$ from Section~\ref{section:Computability_Theory}. The universality of $\mathbf{U}$ implies some invariance and minimality properties in the following senses. Let $\sigma$ and $\tau$ be finite binary strings and $M$ to be an oracle Turing machine on finite binary strings. By universality, it holds that there exists a string $\nu_M$ which may be prepended to the input $\sigma$ to allow $\mathbf{U}$ to simulate $M$ on $\sigma$ given any oracle $B$, i.e., $\mathbf{U}^B(\nu_M\ ^\smallfrown \sigma) = M^B(\sigma)$. Define the \textit{conditional plain Kolmogorov complexity of $\sigma$ given $\tau$ in $M$} to be
\begin{align*}
    C_M(\sigma \mid \tau) = \min \set{\operatorname{len}(\pi)}{M^{\overline{\tau}}(\pi) \downarrow= \sigma},
\end{align*}
or $+\infty$ when minimizing over an empty set; where we map the condition $\tau$ into a prefix-free set of strings via $\overline{\tau} = \tau_0 \tau_0 \tau_1 \tau_1 \cdots \tau_{\operatorname{len}(\tau) - 1} \tau_{\operatorname{len}(\tau) - 1} 01$ (on account of the fact that the set of oracles on which a machine $M$ may converge on a fixed input $\sigma$ must be prefix-free; see \cite{DH10} for more details). Then, call $C(\sigma \mid \tau) := C_{\mathbf{U}}(\sigma \mid \tau)$ the \textit{conditional plain Kolmogorov complexity of $\sigma$ given $\tau$}. 
Under any other choice of universal machine $\mathbf{U}$, this quantity is invariant up to an additive constant independent of $\sigma$ and $\tau$. Conditional plain complexity also satisfies minimality, in that for any $M$:
\begin{align*}
    C(\sigma \mid \tau) \leq C_M(\sigma \mid \tau) + O_M(1).
\end{align*}
for some constant term depending on $M$. The \textit{plain Kolmogorov complexity} of a string $\sigma$ is simply $C(\sigma) = C(\sigma \mid \langle \rangle)$. We call any finitary object $\pi$ satisfying $M^{\overline{\tau}}(\pi) \downarrow = \sigma$ an \textit{$M$-description} for ``$\sigma \mid \tau$''. Any $M$-description $\pi$ of $\sigma \mid \tau$ for which $\operatorname{len}(\pi) = C_M(\sigma \mid \tau)$ is also called \textit{minimal}.

% Add to our vocab a "description" or "program", minimal description, Kolmogorov complexity

A \textit{prefix-free (PF) oracle machine} is an oracle machine on finite binary strings with prefix-free domain for any oracle. Any PF oracle machine $M$ computes a prefix-free partial function $\Psi_M :\subseteq 2^{<\omega} \to 2^{<\omega}$ as before. Again, by the Enumeration Theorem, there exists a \textit{universal} PF oracle machine $\mathbf{U}_{\operatorname{PF}}$. That is, for each PF oracle machine $M$, there exists a coding string $\nu_M \in 2^{<\omega}$ such that given any oracle $B$ and input $\sigma \in 2^{<\omega}$, we have $\mathbf{U}_{\operatorname{PF}}^B(\nu_M\ ^\smallfrown \sigma) = M^B(\sigma)$. Once again, $\mathbf{U}_{\operatorname{PF}}$ admits an invariant and minimal complexity notion in the sense that for any other PF oracle machine $M$, there exists an additive constant term $O_M(1)$ such that for any $\sigma, \tau \in 2^{<\omega}$,
\begin{align*}
    C_{\mathbf{U}_{\operatorname{PF}}}(\sigma \mid \tau) \leq C_M(\sigma \mid \tau) + O_M(1).
\end{align*}
We let $K(\sigma \mid \tau) = C_{\mathbf{U}_{\operatorname{PF}}}(\sigma \mid \tau)$ denote the \textit{conditional prefix Kolmogorov complexity of $\sigma$ given $\tau$}.  Both plain and prefix complexities may be relativized to any oracle by pairing the oracle and the condition string (see \cite{DH10} for more details). Then $K^B(\sigma \mid \tau) := C^B_{\mathbf{U}_{\operatorname{PF}}}(\sigma \mid \tau)$ denotes the \textit{conditional prefix Kolmogorov complexity of $\sigma$ given $\tau$ relative to $B$}. 

These complexity notions may also be extended to any other standard collection of finitary objects $\mathscr{N}$ with a constant amount of overhead, meaning we would only need to sacrifice at most a constant number of bits to describe the computable bijection between $2^{<\omega}$ and $\mathscr{N}$.

\subsection{Results on Prefix Complexity}

We review some standard results on prefix complexity.

% Chain Rule/SOI/Levin-Gács Formula for K(pair)
First, one might refer to the following result as either the \textit{symmetry of information} or the \textit{chain rule} for conditional prefix Kolmogorov complexity. It was first shown by P. G{\'a}cs \cite{Gac74}. The first equation is the strong form and is a property specific to prefix complexity; whereas the second equation is a weak form (where equality holds up to logarithmic terms in the string lengths) shared by many Kolmogorov complexity notions. 

\begin{theorem}[Chain Rule for Conditional Prefix Kolmogorov Complexity; Theorem 1 of \cite{Gac74}]
    \label{theorem:chain_rule_for_K}
    There exists a constant $c \in \omega$ such that for any $\sigma, \tau, \zeta \in 2^{<\omega}$, 
    \begin{align*}
        \Kc{\sigma, \tau}{\zeta} &= \Kc{\sigma}{\zeta} + \Kc{\tau}{\sigma, \K{\sigma}, \zeta} \pm c\\
        &= \Kc{\sigma}{\zeta} + \Kc{\tau}{\sigma, \zeta} \pm [O(\log \operatorname{len}(\sigma)) + c].
    \end{align*}
\end{theorem}

An easy consequence of the chain rule is the \textit{subadditivity} of $K$: i.e., for any $\sigma, \tau \in 2^{<\omega}$,
\begin{align*}
    K(\sigma^\smallfrown \tau) \leq K(\sigma, \tau) \leq K(\sigma) + K(\tau) + O(1).
\end{align*}

% Known bounds on K
We may also say something about the growth-rate of $K$ as an integer function. Write $\operatorname{len}(n) + \operatorname{len}(\operatorname{len}(n)) + \cdots$ for the sum of nested applications of $\operatorname{len}$, including only those terms which are positive. Let $\log^*(n)$ denote the number of such terms in that sum. Thenm, for any natural number $n \in \omega$, the optimized prefix-codes in equation (3.2) of \cite{LV08} demonstrate
\begin{align}
    \K{n} &\leq \log^* n + \operatorname{len}(n) + \operatorname{len}(n)) + \operatorname{len}(\operatorname{len}(\operatorname{len}(n))) + \cdots + O(1) \label{equation:K_bound_for_naturals}\\
    &\leq \log(1 + n) + 2 \log \log (2 + n) + O(1) \nonumber\\
    &= O(\log n) + O(1) \nonumber.
\end{align}
Intuitively, we consider $n$ to be a finite binary string of length $\operatorname{len}(n)$. One description of $n$ involves storing an index (of length no more than $\operatorname{len}(n)$) of where $n$ appears among all other strings of length $\operatorname{len}(n)$ when ordered lexicographically. This algorithm requires a description for $\operatorname{len}(n)$ to work, so we might iterate the argument on $\operatorname{len}(n)$ as appearing somewhere in the list of all strings of length $\operatorname{len}(\operatorname{len}(n))$, and so on. The $\log^* n$ term is necessary because our machine must be prefix-free. 

Together, the chain rule~\ref{theorem:chain_rule_for_K} and the inequality in~\eqref{equation:K_bound_for_naturals} imply the results of Example 3.1.6 of \cite{LV08}. Namely, that for any finitary $n \in \omega$, 
\begin{align}
    \K{n} &\leq \Kc{n}{\operatorname{len}(n)} + \K{\operatorname{len}(n)} + O(1) \nonumber \\
    &\leq \Kc{n}{\operatorname{len}(n)} + \log^* \operatorname{len}(n) + \operatorname{len}(\operatorname{len}(n))+ \cdots + O(1) \label{upper_limit_to_K}\\
    &\leq \Kc{n}{\operatorname{len}(n)} + \log \p{1 + \operatorname{len}(n)} + 2 \log \log (2 + \operatorname{len}(n)) + O(1) \nonumber \\
    &= \Kc{n}{\operatorname{len}(n)} + O(\log \operatorname{len}(n)) + O(1). \nonumber
\end{align}
This fact extends to a complexity bound for any integer tuple based on its Euclidean length (see Observation A.3 of \cite{LL18}): for any tuple-length $m \in \omega$ and integer tuple $z \in \mathbb{Z}^m$,
\begin{align}
    \label{equation:K_bound_int_tuple}
    K(z) &\leq m \cdot \log \p{1 + \norm{z}} + 2 \cdot \log \log \p{2 + \norm{z}} + O_m(1).
\end{align}
The form of this inequality extends to any $r$-dyadic rational tuple $q \in \mathbb{D}^m_r$, where the constant term may depend on both $m$ and $r$. 

% Downey-Hirschfeldt result on K invariant wrt computable functions, maybe even continuity
We include here the conditional version of another standard observation for Kolmogorov complexity: that the complexity of the output of an algorithm is bounded from above by the the complexity of the input plus some overhead to describe the algorithm. This follows from the universality of $\mathbf{U}_{\operatorname{PF}}$. 

\begin{theorem}[Conditional version of Proposition 3.5.4 of \cite{DH10}]
    \label{theorem:bounded_computable_output}
    Let $\Phi : \subseteq 2^{<\omega} \to 2^{<\omega}$ be partial computable. Then, there exists a constant $c \in \omega$ such that for any $\sigma \in \dom(\Phi)$ and $\tau \in 2^{<\omega}$,
    \begin{align*}
        K(\Phi(\sigma) \mid \tau) \leq K(\sigma \mid \tau) + c.
    \end{align*}
\end{theorem}

% Continuity
As a consequence, as noted on page 218 of \cite{LV08}, one may conclude that $K$ is continuous as an integer function. That is, for any natural numbers $n, \Delta n \in \omega$,
\begin{align*}
    \abs{\K{n + \Delta n} - \K{n}} \leq \K{\Delta n} + O(1).
\end{align*}

We might extend this and Theorem~\ref{theorem:bounded_computable_output} as follows. Define $\Phi : \subseteq 2^{<\omega} \to 2^{<\omega}$ to be \textit{partial co-computable (p.c.c.)}  if $\Phi$ is injective and has a partial computable inverse defined on the range of $\Phi$. And $\Phi$ is called \textit{partial bi-computable (p.b.c.)} if $\Phi$ is both partial computable and partial co-computable. 

\begin{proposition}
    \label{proposition:bi_computable_K}
    Let $\Phi : \subseteq 2^{<\omega} \to 2^{<\omega}$. Then there exists a constant $c \in \omega$ such that for any $\sigma \in \dom(\Phi)$ and $\tau \in 2^{<\omega}$,
    \begin{align*}
        \Phi \text{ is p.c.} &\implies K(\Phi(\sigma) \mid \tau) \leq K(\sigma \mid \tau) + c \quad \text{and} \quad K(\tau \mid \sigma) \leq K(\tau \mid \Phi(\sigma)) + c,\\
        \Phi \text{ is p.c.c.} &\implies K(\sigma \mid \tau) \leq K(\Phi(\sigma) \mid \tau) + c \quad \text{and} \quad K(\tau \mid \Phi(\sigma)) \leq K(\tau \mid \sigma) + c,\\
        \Phi \text{ is p.b.c.} &\implies K(\Phi(\sigma) \mid \tau) = K(\sigma \mid \tau) \pm c \quad \text{and} \quad K(\tau \mid \Phi(\sigma)) = K(\tau \mid \sigma) \pm c.
    \end{align*}
\end{proposition}

% Recipe for proving certain inequalities for K by appealing to Church-Turing thesis
Let us describe a general scheme for proving bounds for prefix complexity. Suppose we would like to prove a bound of the form:
\begin{align*}
    K(a \mid b) \leq K(a_1 \mid b_1) + \cdots + K(a_k \mid b_k) + O(1),
\end{align*}
where $a, b$ and $a_1, b_1, ..., a_k, b_k$ are some finitary arguments and $O(1)$ may depend on $k$. In plainer language, this would read as ``one may produce a description for $a$ given $b$ and given descriptions for each $a_i$ given $b_i$.'' 

Formally, to prove a bound of this form, one should construct an appropriate PF oracle machine exhibiting it. The machine $M$ should interpret its input as a tuple of $k$ strings: $\langle \pi_{\widehat{a}_1 \mid \widehat{b}_1}, ..., \pi_{\widehat{a}_k \mid \widehat{b}_k} \rangle$, and read a string $\widehat{b}$ from its oracle tape. The machine may use $\widehat{b}$ and any other results stored along the course of the computation to produce some strings called $\widehat{b}_i$ for any of the $1 \leq i \leq k$. Also, for any $1 \leq i \leq k$ for which $\widehat{b}_i$ is defined, the machine may also simulate $\mathbf{U}_{\operatorname{PF}}$ with oracle $\widehat{b}_i$ and input $\pi_{\widehat{a}_i \mid \widehat{b}_i}$ to compute and report a string called $\widehat{a}_i$ if the computation converges. Of course, it depends on the specific argument how exactly the machine will operate. 

With this oracle machine now defined, one would pass into $M$ the oracle $b$ and the input consisting of minimal $\mathbf{U}_{\operatorname{PF}}$-descriptions $\pi_{a_i \mid b_i}$ for $a_i$ given $b_i$ for each $1 \leq i \leq k$. Supposing that $M$ was properly defined so as to halt on this choice of oracle and input and to correctly report $a$, we may then conclude:
\begin{align*}
    K(a \mid b) &\leq C^b_M(a_1, ..., a_k \mid b_1, ..., b_k, k) + O(1)\\
    &\leq \operatorname{len}(\pi_{a_1 \mid b_1}) + \cdots + \operatorname{len}(\pi_{a_k \mid b_k}) + K(k) + O(1)\\
    &= K(a_1 \mid b_1) + \cdots + K(a_k \mid b_k) + K(k) + O(1).
\end{align*}
It is standard practice to instead verbally explain an (intuitive) algorithm which produces $a$ from $b$ and the other descriptions $a_i \mid b_i$, and then to appeal to the Church-Turing thesis to conclude there exists a prefix-free oracle machine simulating that process. Several of the proofs provided here follow this form for the sake of clarity. 

% Two-part descriptions
One may use this scheme to prove Proposition~\ref{proposition:bi_computable_K}. It also works for \textit{two-part descriptions}. Let $\mathscr{N}$ be any standard collection of finitary objects and $n \in \mathscr{N}$. Given $n \in A$ for some finite set $A \subset \mathscr{N}$, we have that $K(n) \leq K(A) + \log \abs{A} + O(1)$. That is, a description (of a computable enumeration) of $A$, together with a description for the index $i$ when $n$ appears in that enumeration of $A$, suffice to describe $n$ as the $i$-th element enumerated into $A$. In this sense, the pair $(A, i)$ forms a two-part description for $n$. Two-part descriptions are also important in the study of algorithmic statistics \cite{VS16}.

% Kraft Inequality
An easy consequence of $\mathbf{U}_{\operatorname{PF}}$ being prefix-free is the \textit{Kraft Inequality}.
\begin{theorem}[Kraft Inequality, Proposition 3.7.1 of \cite{DH10}]
    \label{theorem:kraft_inequality}
    \begin{align*}
        \sum_{\sigma \in 2^{<\omega}} 2^{-K(\sigma)} \leq 1. 
    \end{align*}
\end{theorem}

% Counting Theorem 
Next, Chaitin's \textit{Counting Theorem} offers useful bounds for the behavior of $K$ on the strings of some fixed length. This essentially follows from the universality of $\mathbf{U}_{\operatorname{PF}}$. 

\begin{theorem}[Chaitin's Counting Theorem \cite{Cha76}, Theorem 3.7.6 of \cite{DH10}]
    \label{theorem:chaitin_counting_theorem}
    There exists a constant $c \in \omega$ such that for any $n, r \in \omega$, 
    \begin{enumerate}[label=(\roman*)]
        \item $\max \set{K(\sigma)}{\sigma \in 2^n} = n + K(n) \pm c$, and 
        \item $\abs{\set{\sigma \in 2^{n}}{K(\sigma) \leq n  + K(n) - r}} \leq 2^{n - r + c}$.
    \end{enumerate}
\end{theorem}

% Note K has an intimate relationship with Shannon Entropy H, including satisfying all the same inequalities (Shen, et al.)
It is interesting to note that the theory of prefix complexity (an algorithmic entropy notion) mimics that of Shannon entropy \cite{Kol68}, and that $K$ satisfies the same linear inequalities true for all inputs as Shannon entropy $H$ up to a logarithmic term in the sum of the lengths of the strings involved. This may be made precise by translating between random variables with finite range and finite binary strings, as is done in \cite{HRSV00}.

\section{Lifting Prefix Complexity to Reals} \label{section:Lifting_Prefix_Complexity_to_Reals}
% Briefly mention Cantor space
Prefix complexity may be lifted to infinitary objects via finitary approximations. We state the following results how they were originally presented: over Euclidean space. Analogous statements also hold over Cantor space. 

\subsection{Lifting Conditional Complexity}

Take any ambient dimensions $m, n \in \omega$, points $x \in \mathbb{R}^m$ and $y \in \mathbb{R}^n$, and precision-levels $r, s \in \omega$. A straightforward definition of the complexity of $x$ given $y$ up to these precision-levels would be to evaluate the complexity: $K(\restrict{x}{r} \mid \restrict{y}{s})$. 

An alternative lift comes from the intuition offered by A.~Shen and N.~Vereshchagin in \cite{SV02}: that for rationals $p \in \mathbb{Q}^m$ and $q \in \mathbb{Q}^n$, the conditional complexity $K(p \mid q)$ could be understood as the complexity of the problem $Q \to P$, where $P$ is the problem ``construct $p$'' and $Q$ is the problem ``construct $q$.'' So, the lift of conditional complexity to arbitrary subsets $X \subseteq \mathbb{R}^m$ given $Y \subseteq \mathbb{R}^n$ should reflect the complexity of constructing a point $p \in X \cap \mathbb{Q}^m$ from any point $q \in Y \cap \mathbb{Q}^n$. This definition was made explicit by J.~Lutz and N.~Lutz in \cite{LL18}.

% Lifts to Euclidean
In particular, they defined the \textit{conditional prefix complexity of $X$ given $Y$} to be:
\begin{align*} 
    \Kc{X}{Y} := \max_{q} \set{\min_{p} \set{\Kc{p}{q}}{p \in X \cap \mathbb{Q}^m}}{q \in Y \cap \mathbb{Q}^n},
\end{align*}
or $+\infty$ if either set is empty. Notice that this notion of complexity is governed only by the rational elements of the sets involved. Unconditionally, the \textit{prefix complexity of $X$} is:
\begin{align*}
    \K{X} := \min\set{\K{p}}{p \in X \cap \mathbb{Q}^m},
\end{align*}
or $+\infty$ whenever $X$ contains no rationals. Thus, a set is as simple as its simplest, rational element. 

Now, for any $x \in \mathbb{R}^m$ and $y \in \mathbb{R}^n$, the \textit{conditional prefix complexity of $x$ to precision-level $r$ given $y$ to precision-level $s$} is:
\begin{align*}
    \Kc[r \mid s]{x}{y} := \Kc{B_{2^{-r}}(x)}{B_{2^{-s}}(y)},
\end{align*}
while the \textit{prefix complexity of $x$ to precision-level $r$ (possibly given $q \in \mathbb{Q}^n$)} is:
\begin{align*}
    K_r(x \mid q) &:= K(B_{2^{-r}}(x) \mid \seq{q}), \text{ and}\\
    K_r(x) &:= K(B_{2^{-r}}(x)).
\end{align*}

In the coming subsection, we will compare $K_{r \mid s}(x \mid y)$ to $K(\restrict{x}{r} \mid \restrict{y}{s})$.

It will also serve us to identify rational elements of a subset which witness the set's prefix complexity. We borrow the nomenclature of A~Case. and J.~Lutz \cite{CL14}.
\begin{definition}
    Let $X \subseteq \mathbb{R}^m$, and suppose $x \in X \cap \mathbb{Q}^m$ satisfies $K(x) \leq K(X) + \varepsilon$ for some $\varepsilon \geq 0$. We say that $x$ is an \textit{$\varepsilon$-approximate $K$-minimizer} of $X$. And if this holds for $\varepsilon = 0$, then $x$ is a \textit{$K$-minimizer} of $X$.
\end{definition}

\subsection{Results on Lifted Conditional Complexity}

We recall two lemmas about $K_{r \mid s}$ known as its \textit{linear sensitivities}. 

\begin{lemma}[Linear Sensitivity of $K_r$ in $r$; Lemma 3.8 of \cite{CL14}]
    For all $m \in \omega$, there exists $c \in \omega$ such that for all $x \in \mathbb{R}^m$ and $r, \Delta r \in \omega$,
    \begin{align*}
        K_r(x) \leq K_{r + \Delta r}(x) \leq K_r(x) + K(r) + m \Delta r + a_{\Delta r} + c,
    \end{align*}
    where $a_{\Delta r} := K(\Delta r) + 2 \log \p{\ceil{\frac{1}{2} \log m} + \Delta r + 3} + \p{\ceil{\frac{1}{2} \log m} + 3} m + K(m) + 2 \log m$.
    \label{lemma:linearity_r_no_condition}
\end{lemma}

\begin{lemma}[Linear Sensitivity of $K_{r \mid s}$ in $s$; Lemmas 7 and 8 of \cite{LL18}]
    \label{lemma:linearity_with_condition}
    For all $m, n \in \omega$, there exists $c \in \omega$ such that for all $x \in \mathbb{R}^m$, $y \in \mathbb{R}^n$, $q \in \mathbb{Q}^n$, and $r, s, \Delta r, \Delta s \in \omega$,
    \begin{enumerate}[label=(\roman*)]
        \item $K_r(x \mid q) \leq K_{r + \Delta r}(x \mid q) \leq K_{r}(x \mid q) + m \Delta r + 2 \log(1 + \Delta r) + K(r, \Delta r) + c$;
        \item $K_{r \mid s}(x \mid y) \geq K_{r \mid s + \Delta s}(x \mid y) \geq K_{r \mid s}(x \mid y) - n \Delta s - 2 \log(1 + \Delta s) + K(s, \Delta s) + c$.
    \end{enumerate}
\end{lemma}

A consequence of the linear sensitivities is that one may approximate the conditional prefix complexity of Euclidean points more simply via dyadic truncations.

\begin{lemma}[Lemma~A.1 of \cite{LS16}]
    \label{maxmin_truncation_for_r}
    For all $m, n \in \omega$, there exists $c \in \omega$ such that for all $x \in \mathbb{R}^m$, $q \in \mathbb{Q}^n$, and $r \in \omega$,
    \begin{align*}
        \abs{\Kc[r]{x}{q} - \Kc{\restrict{x}{r}}{q}} \leq \K{r} + c.
    \end{align*}
\end{lemma}

\begin{lemma}[Lemma~A.3 of \cite{LS16}]
    \label{maxmin_truncation_for_r_s}
    For all $m, n \in \omega$, there exists $c \in \omega$ such that for all $x \in \mathbb{R}^m$, $y \in \mathbb{R}^n$, and $r, s \in \omega$,
    \begin{align*}
        \abs{\Kc[r \mid s]{x}{y} - \Kc{\restrict{x}{r}}{\restrict{y}{s}}} \leq \K{r} + \K{s} + c.
    \end{align*}
\end{lemma}

One may then conclude an approximate symmetry of information for prefix complexity on Euclidean space. Again, this is a weak form of symmetry of information.

\begin{theorem}[Approximate Chain Rule; Lemma~4 of \cite{LS16}]
    \label{theorem:approximate_SOI_conditional}
    For any $m, n \in \omega$, $x \in \mathbb{R}^m$, $y \in \mathbb{R}^n$, and $r \geq s \in \omega$,
    \begin{enumerate}[label=(\roman*)]
        \item $\abs{K_{r \mid r}(x \mid y) + K_r(y) - K_{r}(x,  y)} \leq O_{m, n}(\log r) + O_n(\log \log \norm{y}) + O_{m, n}(1)$;
        \item $\abs{K_{r \mid s}(x \mid x) + K_s(x) - K_r(x)} \leq O_m(\log r) + O_m(\log \log \norm{x}) + O_m(1)$.
    \end{enumerate}
\end{theorem}

We note that when computing with infinitary objects such as infinite binary strings or Euclidean reals, one must distinguish between \textit{oracle} and \textit{conditional} accesses. While oracle access permits computations with approximations of the given data to arbitrary precisions, conditional access places a firm limit on the precision-level to which one may access the given data.

\begin{lemma}[Lemma~14 of \cite{LL18}]
    \label{oracle_access_to_conditional}
    For all $m, n \in \omega$, there exists a constant $c \in \omega$ such that, for all $x \in \mathbb{R}^m$, $y \in \mathbb{R}^n$, and $r, s \in \omega$,
    \begin{align*}
        K^y_r(x) \leq K_{r \mid s}(x \mid y) + K(s) + c.
    \end{align*}
\end{lemma}

\section{Semimeasures} \label{section:Semimeasures}

This section presents an alternative approach to measuring the information-content or algorithmic-complexity of finitary objects through mass distribution. 

\subsection{Discrete Semimeasures}
% Define discrete semimeasures, lower semicomputability
A function $m: 2^{<\omega} \to [0, 1]$ is said to be a \textit{discrete semimeasure} on $2^\omega$ if it satisfies $\sum_{\sigma \in 2^{<\omega}} m(i) \leq 1$. There exists a (reference) \textit{optimal} lower-semicomputable discrete semimeasure $\mathbf{m}$ on $2^\omega$. That is, if $m$ is any other lower-semicomputable discrete semimeasure, there exists a constant $\beta > 0$ such that for all $\sigma \in 2^{<\omega}$, we have $\mathbf{m}(\sigma) \geq \beta \cdot m(\sigma)$. Refer to Section 4.3 of \cite{LV08} for more details. 

% Coding Theorem
There is a natural connection between prefix complexity and lower-semicomputable discrete semimeasures. Levin originally introduced the concept of an \textit{information content measure}, or an upper-semicomputable, partial function $F : \subseteq 2^{<\omega} \to \omega$ satisfying $\sum_{\sigma \in 2^{<\omega}} 2^{-F(\sigma)} \leq 1$, as described in \cite{DH10}. By the Kraft Inequality (\ref{theorem:kraft_inequality}) and the universality of $\mathbf{U}_{\operatorname{PF}}$, $K$ is considered to be minimal among the the information content measures. Similarly, $2^{-K(\cdot)}$ is maximal (i.e., optimal) among the lower-semicomputable discrete semimeasures.

This relationship is made more precise in \textit{Levin's Coding Theorem}. We start by reviewing two classes of lower-semicomputable discrete semimeasures on $2^{<\omega}$. 

\begin{definition}
    Fix an oracle machine $M$ on $2^{<\omega}$, and $\sigma \in 2^{<\omega}$. 
    \begin{itemize}
        \item Let $R_M(\sigma) := 2^{-C_M(\sigma)}$ denote the \textit{algorithmic probability} of $\sigma$ under $M$.
        \item Let $Q_M(\sigma) := \sum_{M(\pi) \downarrow = \sigma} 2^{- \operatorname{len}(\pi)}$ denote the \textit{a priori probability} of $\sigma$ under $M$.
    \end{itemize}
    When $M$ is the reference universal PF oracle machine, simply write $R \equiv R_{\mathbf{U}_{\operatorname{PF}}}$ and $T \equiv T_{\mathbf{U}_{\operatorname{PF}}}$. 
\end{definition}

Levin proved that $\mathbf{m}$, $R$, and $Q$ all capture essentially the same quantity.

\begin{theorem}[Levin's Coding Theorem, Theorem~3.9.4 of \cite{DH10}]
    \label{theorem:coding_theorem}
    For any $\sigma \in 2^{<\omega}$,
    \begin{align*}
        K(\sigma) = -\log R(\sigma) = -\log Q(\sigma) \pm O(1) = -\log \mathbf{m}(\sigma) \pm O(1).
    \end{align*} 
\end{theorem}

One informal corollary of the Coding Theorem is that: having many \textit{long} $\mathbf{U}_{\operatorname{PF}}$-descriptions implies having at least one \textit{short} $\mathbf{U}_{\operatorname{PF}}$-description. To see this, suppose $\Pi$ is a collection of strings $\pi \in 2^{<\omega}$ all of a common length $\ell \in \omega$ and serving as $\mathbf{U}_{\operatorname{PF}}$-descriptions for the string $\sigma \in 2^L$. Then, the \textit{a priori} probability of $\sigma$ will be at least:
\begin{align*}
    Q(\sigma) &= \lambda([\set{\pi \in 2^{< \omega}}{\mathbf{U}_{\operatorname{PF}}(\pi) \downarrow = \sigma}]) \geq \lambda([\Pi]) = \sum_{\pi \in \Pi} 2^{-\operatorname{len}(\pi)} = \abs{\Pi} \cdot 2^{-\ell}. 
\end{align*}
Thus, by the Coding Theorem~\ref{theorem:coding_theorem}, $\sigma$ has an appropriately shorter $\mathbf{U}_{\operatorname{PF}}$-description:
\begin{align*}
    K(\sigma) \leq - \log Q(\sigma) + O(1) \leq \ell - \log \abs{\Pi} + O(1). 
\end{align*}

\subsection{Continuous Semimeasures}
Discrete semimeasures do not take into account the topology of Cantor space, while a continuous semimeasure does by considering how mass distributes from a string to its extensions. In particular, a function $M: 2^{<\omega} \to [0, 1]$ is said to be a \textit{continuous semimeasure} on $2^\omega$ if it satisfies $M(\sigma) \geq M(\sigma^\smallfrown 0) + M(\sigma ^\smallfrown 1)$ for all $\sigma \in 2^{<\omega}$. As in the discrete case, there exists a reference \textit{optimal} lower-semicomputable continuous semimeasure $\mathbf{M}$ on $2^\omega$. That is, if $M$ is any other lower-semicomputable continuous semimeasure, there exists a constant $\beta > 0$ such that for all $\sigma \in 2^{<\omega}$, we have $\mathbf{M}(\sigma) \geq \beta \cdot M(\sigma)$. 
% This $\mathbf{M}$ plays the role of a \textit{universal prior} in the context of Solomonoff induction. 
Refer to Section 4.5 of \cite{LV08} for more details. 
% and 5.2 for Solomonoff induction

% A priori complexity
We denote by $\operatorname{KM}(\sigma) := - \log \mathbf{M}(\sigma)$ the \textit{a priori complexity} of a string $\sigma \in 2^{<\omega}$. It is interesting to compare this complexity notion to prefix complexity. There are few satisfying bounds on their difference. Refer to \cite{US96} and Section 4.5.5 of \cite{DH10} for more details.

It is important to mention the fundamental role that constructive martingales have played in the development of algorithmic information theory. By 2000, J.~Lutz had introduced algorithmically-restricted betting strategies called \textit{constructive supergales} as a generalization to the constructive martingales in C. Schnorr's early work on algorithmic randomness \cite{Lut00, Sch71}. Let us see why, for a certain class of premeasures, these betting strategies are interchangeable with the lower-semicomputable continuous semimeasures. 

Given a premeasure $\rho$ on $2^\omega$, a \textit{$\rho$-supergale} is a map $\delta : 2^{<\omega} \to [0, +\infty)$ such that
\begin{align*}
    \delta(\sigma) \cdot \rho([\sigma]) \geq \delta(\sigma^\smallfrown 0) \cdot \rho([\sigma^\smallfrown 0]) + \delta(\sigma^\smallfrown 1) \cdot \rho([\sigma^\smallfrown 1]),
\end{align*}
for each $\sigma \in 2^{<\omega}$. If indeed $\rho = \rho_{s}$ for some $s \geq 0$, simply refer to $\delta$ as an \textit{$s$-supergale}.

A $\rho$-supergale is said to \textit{succeed} on a subset $X \subseteq 2^{\omega}$ if for each $x \in X$, we have $\limsup_{r \to \infty} \delta(\restrict{x}{r}) = + \infty$. Thus, to succeed on an infinite sequence $x \in 2^\omega$ is to earn arbitrarily much by taking bets on new bits of $x$ following the betting strategy $\delta$. 

It had already been shown by J.~Lutz in \cite{Lut2000} that the Hausdorff dimension of a subset $X \subseteq 2^\omega$ may be characterized by the ability for $s$-supergales to succeed on $X$. That is,
\begin{align}
    \label{equation:classical_relationship_supergales_and_Hausdorff_dimension}
    \dim_{\operatorname{H}}(X) = \inf \set{s \geq 0}{\text{there exists an $s$-supergale succeeding on $X$}}.
\end{align}
Lutz's subsequent idea was to restrict supergales by an effectivity condition and see whether they could still recover classical Hausdorff dimension. Call a $\rho$-supergale to be \textit{constructive} if it is lower-semicomputable. If $\rho$ is length-invariant, strictly-positive, and computable on the cylinder sets $\mathcal{B}$ (e.g., any $s$-dimensional Hausdorff premeasure $\rho_{s}$ with computable $s \geq 0$), it holds that $\delta$ is a (constructive) $\rho$-supergale if and only if the map $M : \sigma \mapsto \delta(\sigma) \cdot \rho([\sigma])$ is a (lower-semicomputable) continuous semimeasure. 

Thus, continuous semimeasures are also related to Hausdorff dimension in $2^\omega$, and we may extend the notion of success to them as follows: a continuous semimeasure $M$ is said to \textit{$\rho$-succeed} on $X \subseteq 2^\omega$ if for each $x \in X$, we have $\limsup_{r \to \infty} M(\restrict{x}{r}) / \rho([\restrict{x}{r}]) = +\infty$. If $M$ is some lower-semicomputable continuous semimeasure which $\rho$-succeeds on $X$, then, by optimality, $\mathbf{M}$ also $\rho$-succeeds on $X$. When $\rho = \rho_{s}$, we refer to this as \textit{$s$-success}. 

From this perspective, the language of constructive supergales ports into the language of lower-semicomputable continuous semimeasures. 

% Correspondence between continuous semimeasures and strong Solovay tests
Another characterization of $s$-success for lower-semicomputable continuous semimeasures has to do with the existence of strong Solovay-$s$-tests of arbitrarily small prefix $s$-weight. Call $(q_n)_{n \in \omega} \subseteq \mathbb{Q}_{>0}$ \textit{rapidly decreasing} if $\lim_{n \to \infty} q_n \cdot 2^n = 0$.

\begin{theorem}
    \label{theorem:continuous_semimeasure_is_supergale_is_strong_Solovay_tests}
    Fix $X \subseteq 2^\omega$ and computable $s > 0$. Then the following are equivalent.
    \begin{enumerate}[label = (\roman*)]
        \item There is a uniformly-c.e.~sequence of strong Solovay-$s$-tests $(\mathcal{V}_n)_n$ and a computable, rapidly decreasing sequence $(q_n)_n$ where each $\mathcal{V}_n$ covers $X$ and satisfies $\operatorname{PW}_s(\mathcal{V}_n) \leq q_n$; 
        \item Some constructive $s$-supergale succeeds on $X$;
        \item Some lower-semicomputable continuous semimeasure $s$-succeeds on $X$; and
        \item $\mathbf{M}$ $s$-succeeds on $X$.
    \end{enumerate}
\end{theorem}
\begin{proof}
    It suffices to prove (i) $\iff$ (iii). First, suppose (i). For each $n \in \omega$, the map
    \begin{align*}
        M_n(\sigma) := \frac{1}{q_n} \operatorname{PW}_s(\set{\tau\in \mathcal{V}_n}{\sigma \preceq \tau})
    \end{align*}
    is a lower-semicomputable continuous semimeasure. And, whenever $\sigma \in \mathcal{V}_n$, it holds that $M_n(\sigma) \geq \frac{1}{q_n} \cdot 2^{-s \cdot \operatorname{len}(\sigma)}$. Now, define another lower-semicomputable continuous semimeasure
    \begin{align*}
        M(\sigma) := \sum_{n \in \omega} \frac{M_n(\sigma)}{2^{n + 1}}.
    \end{align*}
    Then, for each $\sigma \in \mathcal{V}_n$, we have $M(\sigma) \geq \frac{1}{q_n} \cdot 2^{-(n + 1) - s \cdot \operatorname{len}(\sigma)}$. Now, if $x \in X$ and $X$ is covered by each $\mathcal{V}_n$, it holds that for each $n \in \omega$, there exist infinitely many $r \in \omega$ for which $\restrict{x}{r} \in \mathcal{V}_n$. Hence,
    \begin{align*}
        \limsup_{r \to \infty} \frac{M(\restrict{x}{r})}{2^{-s \cdot \operatorname{len}(\restrict{x}{r})}} = \limsup_{r \to \infty} M(\restrict{x}{r}) \cdot 2^{s \cdot r} \geq \frac{1}{q_n} \cdot 2^{-(n + 1)} \to +\infty
    \end{align*}
    as $n \to +\infty$ since $(q_n)_n$ is rapidly decreasing. Thus, $M$ $s$-succeeds on $X$. 

    In the other direction, suppose (iii). Fix a computable, rapidly decreasing sequence $(q_n)_n$ and define for each $n \in \omega$ the c.e.~set:
    \begin{align*}
        \mathcal{V}_n := \set{\sigma \in 2^{<\omega}}{M(\sigma) \cdot 2^{s \cdot \operatorname{len}(\sigma)} \geq \frac{1}{q_n}}.
    \end{align*}
    Notice that if $\mathcal{P} \subseteq \mathcal{V}_n$ is prefix-free, then
    \begin{align*}
        \operatorname{DW}_s(\mathcal{P}) = \sum_{\sigma \in \mathcal{P}} 2^{- s \cdot \operatorname{len}(\sigma)} \leq q_n \cdot \sum_{\sigma \in \mathcal{P}} M(\sigma) \leq q_n,
    \end{align*}
    so $\operatorname{PW}_s(\mathcal{V}_n) \leq q_n$ for all $n \in \omega$. Finally, for a fixed $n \in \omega$, if $x \in X$, then $M$ $s$-succeeding on $X$ implies there are infinitely many $r \in \omega$ for which $M(\restrict{x}{r}) \cdot 2^{s \cdot r} \geq \frac{1}{q_n}$, or $\restrict{x}{r} \in \mathcal{V}_n$. So, each $\mathcal{V}_n$ is a strong Solovay-$s$-test covering $X$, and the sequence $(\mathcal{V}_n)_n$ is uniformly-c.e.
\end{proof}

\section{Effective Hausdorff Dimension} \label{section:Effective_Hausdorff_Dimension}
J.~Lutz was the first to introduce a refinement to Hausdorff dimension in the form of \textit{constructive dimension} \cite{Lut00}, which could distinguish between subsets of zero Hausdorff dimension by their algorithmically-exploitable patterns. Shortly after, E.~Mayordomo demonstrated this quantity matched an algorithmic information density involving Kolmogorov complexity first studied by L.~Staiger \cite{May02, Sta89}. In this section, we review a few of the characterizations for effective dimension, but not necessarily in their historical order.

% Defined using ML-s-tests
% Existence of a universal ML-s-test
% Stability properties
% Note that a similar method can be used to effectivize packing dimension, refer to Jan's thesis(?), as well as other fractal dimensions like box counting dimension (again in Jan's thesis)

For simplicity, we define the effective dimension notions in this section over $2^\omega$. With slight adjustments, some may be extended to any $\mathbb{R}^m$.

Inspired by the definition of the local dimension of an outer measure in K. Falconer's book \cite{Fal90}, define local dimension for outer measures and semimeasures as follows.

\begin{definition}
    Let $\mu$ either be a discrete semimeasure, continuous semimeasure, or outer measure on $2^\omega$. Then the \textit{(lower) local dimension} of a subset $X \subseteq 2^\omega$ with respect to $\mu$ is defined as:
    \begin{align*}
        \dim_{\operatorname{loc}} \mu(X) = \sup_{x \in X} \liminf_{r \to \infty} \frac{-\log \mu([\restrict{x}{r}])}{r},
    \end{align*}
    where $\mu([\restrict{x}{r}])$ is understood as $\mu(\restrict{x}{r})$ whenever $\mu$ is a semimeasure.
\end{definition}

In the language of supergales, the map $\boldsymbol{\delta}(\sigma) := 2^{s \cdot \operatorname{len}(\sigma)} \cdot \mathbf{M}(\sigma)$ is an \textit{optimal} constructive $s$-supergale whenever $s \geq 0$ is left-c.e. It follows from our Theorem~\ref{theorem:continuous_semimeasure_is_supergale_is_strong_Solovay_tests} that  $\dim_{\operatorname{loc}} \mathbf{M}$ matches J. Lutz's \textit{constructive dimension} \cite{Lut00, Lut03}. That is, for any $X \subseteq 2^\omega$,
\begin{align*}
    \dim_{\operatorname{loc}} \mathbf{M}(X) &= \inf \set{s \geq 0}{\text{there exists a constructive $s$-supergale succeeding on $X$}}\\
    &=: \operatorname{cdim}(X).
\end{align*}
In light of~\eqref{equation:classical_relationship_supergales_and_Hausdorff_dimension}, we may consider $\operatorname{cdim}$ to be an effective version of Hausdorff dimension. We will cite some other asymptotic coincidences further justifying this view. 

In \cite{May02}, Mayordomo proved that Lutz's constructive dimension of a point $x$ matches the \textit{lower incompressibility ratio of prefix Kolmogorov complexity} of $x$: 
\begin{align*}
    \operatorname{cdim}(X) = \underline{\kappa}(X) := \sup_{x \in X} \liminf_{r \to \infty} \frac{K(\restrict{x}{r})}{r} = \dim_{\operatorname{loc}} \mathbf{m}(X),
\end{align*}
where the last equality follows from the Coding Theorem~\ref{theorem:coding_theorem}. So $\mathbf{M}$ and $\mathbf{m}$ produce the same local dimension notion on $2^\omega$. 

Note that by definition, $\dim_{\operatorname{loc}} \mathbf{M}$ is itself the \textit{lower incompressibility ratio with respect to a priori complexity}:
\begin{align*}
    \dim_{\operatorname{loc}} \mathbf{M}(X) = \sup_{x \in X} \liminf_{r \to \infty} \frac{\operatorname{KM}(\restrict{x}{r})}{r}.
\end{align*}

%----
% Effectivized Hausdorff dimension

One might also attempt to effectivize the definition of classical Hausdorff dimension directly. Suppose $\mathcal{U}$ is a Martin-L{\"o}f-$s$-test covering $X \subseteq 2^\omega$. Then, $\mathcal{U}$ acts as an effective witness to $X$ being $s$-null. In analogy with the definition of Hausdorff dimension in~\eqref{equation:definition_of_Hausdorff_dimension}, define the \textit{effective Hausdorff dimension} of $X \in 2^\omega$ to be:
\begin{align*}
    \operatorname{effdim}(X) := \inf \set{s \geq 0}{\text{$X$ is ML-$s$-null}}.
\end{align*}
Note that the same could be done for any other of the effective test notions defined above. 
We prefer to work with Martin-L{\"o}f-randomness because it admits universal tests.
In particular, for each left-c.e.~$s > 0$, there exists a \textit{universal} ML-$s$-test $\mathcal{U}^s$. This means that if $\mathcal{U}$ is an ML-$s$-test covering some subset $X \subseteq 2^\omega$, then so too will $\mathcal{U}^s$ cover $X$. The existence of such a $\mathcal{U}^s$ implies the \textit{pointwise stability} of effective Hausdorff dimension. That is, for any $X \subseteq 2^{\omega}$,
\begin{align*}
    \operatorname{effdim}(X) = \sup_{x \in X} \operatorname{effdim}(\seq{x}).
\end{align*}

Martin-L{\"o}f-$s$-randomness relates to success by constructive $s$-supergales. That is, if there exists a constructive $s$-supergale succeeding on $X$, then $X$ will be ML-$t$-null for all left-c.e.~$t > s$. Conversely, if $s$ is left-c.e.~and $X$ is ML-$s$-null, then there exists a constructive $s$-supergale succeeding on $X$. Therefore, $\operatorname{cdim}(X) = \operatorname{effdim}(X)$ on all sets $X \subseteq 2^\omega$. 

%----
% Locally optimal outer measures

In his doctoral dissertation, N.~Lutz introduced the map $\boldsymbol{\kappa}(X) := 2^{-K(X)}$, an outer measure on $2^\omega$ (originally, $\mathbb{R}^m$) which possesses many algorithmic optimality properties as explored in both \cite{Lut17b, LL20}. In particular, J.~Lutz and N.~Lutz showed that $\boldsymbol{\kappa}(X)$ is \textit{locally optimal} among all the other strongly finite, lower-semicomputable outer measures on $2^\omega$, and that any locally optimal outer measure $\boldsymbol{\mu}$ has local dimension $\dim_{\operatorname{loc}}\boldsymbol{\mu}$ matching the lower incompressibility ratio of prefix complexity \cite{LL20}. 

We now summarize these standard asymptotic coincidences in the following theorem.

\begin{theorem}[Effective Dimension]
    \label{theorem:asymptotic_coincidences_Cantor}
    For any $X \subseteq 2^\omega$, we have
    \begin{align*}
        \small 
        \operatorname{effdim}(X) = \dim_{\operatorname{loc}} \mathbf{M}(X) = \dim_{\operatorname{loc}} \mathbf{m}(X) = \dim_{\operatorname{loc}} \boldsymbol{\kappa}(X) = \operatorname{cdim}(X) = \underline{\kappa}(X).
    \end{align*}
\end{theorem}

We will use $\dim(\cdot)$ as a generic placeholder for any one of those listed in Theorem~\ref{theorem:asymptotic_coincidences_Cantor} over either Cantor space or Euclidean space. And $\dim(x)$ will stand for $\dim(\seq{x})$ for any $x \in 2^\omega$. 

% Note that the upper incompressibility ratio of prefix complexity matches effectivized packing dimension
We further note that \textit{effectivized packing dimension} (or \textit{strong dimension}) coincides with the \textit{upper} incompressibility ratio of prefix complexity, as first shown in  \cite{AHLE07}. So, we fix the notation:
\begin{align*}
    \operatorname{Dim}(X) := \overline{\kappa}(X) = \sup_{x \in X} \limsup_{r \to \infty} \frac{K(\restrict{x}{r})}{r}.
\end{align*}

%--------------

\section{Effective Dimension Variants} \label{section:Conditional_and_Mutual_Dimensions}

We review two variants of effective dimension.

\subsection{Mutual Dimension}

One variant of effective dimension is \textit{mutual dimension}, which was originally defined in \cite{CL14}. It intuitively captures the algorithmic information shared by two infinitary objects in the limit. In Shannon information theory, mutual information is a measure of the correlation between two distributions. A.~Case and J.~Lutz's algorithmic notion of mutual information serves as a refinement of the classical notion. Namely, the \textit{(conditional) mutual information between two finitary objects $a, b \in \mathscr{N}$ given $c \in \mathscr{N}$} is
\begin{align*}
    I(a : b \mid c) := K(a \mid c) - K(a \mid b, c).
\end{align*}
Extending their notation slightly, we define the \textit{mutual information between subsets $X \subseteq \mathbb{R}^m$ and $Y \subseteq \mathbb{R}^n$} as:
\begin{align*}
    I(X : Y) := \min \set{I(p : q)}{p \in X \cap \mathbb{Q}^m \text{ and } q \in Y \cap \mathbb{Q}^n},
\end{align*}
or zero if taken over the empty set. Then the \textit{mutual information between $x \in \mathbb{R}^m$ to precision-level $r$ and $y \in \mathbb{R}^n$ to precision-level $s$} is defined as:
\begin{align*}
    I_{r : s}(x : y) := I(B_{2^{-r}}(x) : B_{2^{-s}}(x)),
\end{align*}
Both $I$ and $I_{r : s}$ inherit many of their expected properties such as a chain rule and symmetry from $K$ and $K_{r \mid s}$. In particular, a slight generalization to Theorem 4.10 in \cite{CL14} for distinct precision-levels proves the chain rule for $I_{r : s}$.
\begin{theorem}[Chain Rule for Mutual Information, c.f. Theorem 4.10 of \cite{CL14}]
    \label{theorem:chain_rule_mutual_information_r_s}
    For all $x \in \mathbb{R}^m$, $y \in \mathbb{R}^n$, and $r, s \in \omega$,
    \begin{align*}
        I_{r : s}(x : y) &= K_r(x) + K_s(y) - K_{r, s}(x, y) \pm (o(r) + o(s))\\
        &= K_r(x) + K_{s \mid r}(y \mid x) \pm (o(r) + o(s)).
    \end{align*}
\end{theorem}

Case and Lutz further defined the \textit{lower and upper mutual dimensions of $x \in \mathbb{R}^m$ and $y \in \mathbb{R}^n$} as follows:
\begin{align*}
    \operatorname{mdim}(x : y) := \liminf_{r \to \infty} \frac{I_{r : r}(x : y)}{r}, \quad \text{and} \quad \operatorname{Mdim}(x : y) := \limsup_{r \to \infty} \frac{I_{r : r}(x : y)}{r}.
\end{align*}
Both mutual dimensions possess a few special properties. For instance, if $x \in \mathbb{R}^m$ and $y \in \mathbb{R}^n$, then $0 \leq \operatorname{mdim}(x : y) \leq \min \seq{\operatorname{mdim}(x), \operatorname{mdim}(y)}$. As well as symmetry: $\operatorname{mdim}(x : y) = \operatorname{mdim}(y : x)$. Analogous statements hold for $\operatorname{Mdim}$. 

\subsection{Conditional Dimension}
A conditional version of effective dimension was first defined in \cite{LL18} using the lifted conditional prefix complexity $K_{r \mid s}$. The \textit{lower and upper conditional dimensions of $x \in \mathbb{R}^m$ given $y \in \mathbb{R}^n$} were defined as:
\begin{align*}
    \dim(x \mid y) = \liminf_{r \to \infty} \frac{K_{r \mid r}(x \mid y)}{r}, \quad \text{and} \quad \operatorname{Dim}(x \mid y) = \limsup_{r \to \infty} \frac{K_{r \mid r}(x \mid y)}{r}.
\end{align*}

Conditional dimension may equivalently be characterized in terms of semimeasures and supergales as in Theorem~\ref{theorem:asymptotic_coincidences_Cantor}. 
% For instance, conditional continuous semimeasures where $\mathbf{M}(\sigma \mid \tau) := \frac{\mathbf{M}(\sigma^\smallfrown \tau)}{\mathbf{M}(\tau)}$ \cite{DH10}.

These conditional dimensions were shown to relate with the weak and strong effective dimensions as follows:
\begin{theorem}[Chain Rule for Conditional Dimension; Corollary 13 of \cite{LL18}]
    \small 
    \begin{align*}
        \dim(x) + \dim(y \mid x) \leq \dim(x, y) &\leq \dim(x) + \operatorname{Dim}(y \mid x) \leq \operatorname{Dim}(x, y) \leq \operatorname{Dim}(x) + \operatorname{Dim}(y \mid x).
    \end{align*}
\end{theorem}

\subsection{Results on Mutual and Conditional Dimensions}

Both mutual dimensions behave predictably under computable, uniformly-continuous maps. We cite here just one corollary of this robustness. 

\begin{theorem}[Preservation of Mutual Dimension; Corollary 8.3 of \cite{CL14}]
    \label{theorem:mutual_dimension_invariant_under_computable_bi_Lipshictz}
    If $f : \mathbb{R}^m \to \mathbb{R}^k$ and $g : \mathbb{R}^n \to \mathbb{R}^\ell$ are computable and bi-Lipschitz, then for all $x \in \mathbb{R}^m$ and $\mathbb{R}^n$,
    \begin{align*}
        \operatorname{mdim}(f(x) : g(y)) &= \operatorname{mdim}(x : y),\\
        \operatorname{Mdim}(f(x) : g(y)) &= \operatorname{Mdim}(x : y).
    \end{align*}
\end{theorem}

In his dissertation \cite{Rei04}, J.~Reimann argues why effective dimension is also preserved under bi-computable, (locally) bi-Lipschitz continuous maps. This fact was used by N.~Lutz to confirm an analog of Marstrand's Projection Theorem for the effective dimension of points \cite{Lut17}. Results like these mimic the preservation of Hausdorff dimension under locally bi-Lipschitz continuous maps (e.g., Corollary 2.4 in \cite{Fal90}). We will expand on this later in Section~\ref{section:Effective_Dimension_under_Uniformly_Continuous_Maps}.

Finally, we summarize the known relations between mutual, conditional, and effective dimensions. 

\begin{proposition}
    \label{proposition:relating_mutual_dimension_and_conditional_dimension}
    For all $x \in \mathbb{R}^m$ and $y \in \mathbb{R}^n$,
    \begin{align*}
        \dim(x \mid y) & \leq \operatorname{dim}(x) - \operatorname{mdim}(x : y) \leq \operatorname{Dim}(x \mid y),\\
        \dim(x \mid y) & \leq \operatorname{Dim}(x) - \operatorname{Mdim}(x : y) \leq \operatorname{Dim}(x \mid y).
    \end{align*}
\end{proposition}
We note that two of these inequalities were already proved in Corollary 11 of \cite{LL18}. Our contribution is to add the uppermost bound in the top line, and the lowermost bound in the bottom line. All of these inequalities follow by re-expressing mutual dimension using the chain rules for $I_{r : r}$ and $K_{r \mid r}$ as in Theorems~\ref{theorem:chain_rule_mutual_information_r_s} and~\ref{theorem:approximate_SOI_conditional}, respectively, and then distributing the limits inferior or superior appropriately over those terms. 

\section{Complexity Characterizations of Algorithmic Randomness} \label{section:Complexity_Characterizations_of_Algorithmic_Randomness}

\subsection{Algorithmic Discrepancies}

% Algorithmic discrepancies 
We fix some notation for the quantitative discrepancy between a finitary object's algorithmic entropy and its mass according to some premeasure. 

\begin{definition}
    Let $\rho$ be a premeasure on $2^\omega$. Define for any $\sigma \in 2^{<\omega}$ its:
    \begin{itemize}
        \item \textit{A priori discrepancy}: $\operatorname{\Delta M}_\rho(\sigma) := \operatorname{KM}(\sigma) + \log \rho([\sigma]) = \log \frac{\rho([\sigma])}{\mathbf{M}(\sigma)}$,
        \item \textit{Prefix discrepancy}: $\operatorname{\Delta P}_\rho(\sigma) := K(\sigma) + \log \rho([\sigma]) =^* \log \frac{\rho([\sigma])}{\mathbf{m}(\sigma)}$.
    \end{itemize}
\end{definition}

When $\rho = \rho_s$, simply write $\operatorname{\Delta M}_s$ and $\operatorname{\Delta P}_s$. 

\subsection{Complexity Characterizations of Algorithmic Randomness}

It turns out that these algorithmic discrepancies characterize the various $s$-dimensional randomness notions defined above. 

\begin{theorem}
    \label{theorem:complexity_characterizations_of_randomness}
    Let $s \geq 0$ be left-c.e.~and $x \in 2^\omega$. Then $x$ is:
    \begin{itemize}
        \item ML-$s$-random $\iff$ $x$ is weak Chaitin-$s$-random: $\operatorname{\Delta P}_s(\restrict{x}{r}) \geq O(1)$, \cite{Tad02};
        \item Solovay-$s$-random $\iff$ $x$ is strong Chaitin-$s$-random: $\operatorname{\Delta P}_s(\restrict{x}{r}) \to \infty$, \cite{Tad02, CST06};
        \item Strong ML-$s$-random $\iff$ $\operatorname{\Delta M}_s(\restrict{x}{r}) \geq O(1)$, \cite{CST06};
        \item Vehement ML-$s$-random $\iff$ $x$ is strong ML-$s$-random;
        \item Strong Solovay-$s$-random $\iff$ $ \operatorname{\Delta M}_s(\restrict{x}{r}) \to \infty$, \cite{MR09}. 
    \end{itemize}
\end{theorem}

In light of these characterizations, one may argue on the level of discrete or continuous semimeasures rather than covers in order to effectivize results on Hausdorff dimension. 

\section{Point-to-Set Principles} \label{section:PTS}

In this context, a point-to-set (PTS) principle will characterize the fractal dimension or measure content of a subset by some local dimension or measure assigned to its elements. We summarize the few PTS principles known for Hausdorff dimension and the $s$-dimensional Hausdorff outer measures via algorithmically-restricted quantities. 

\subsection{PTS for Hausdorff Dimension}

Certain fractal dimension notions have pointwise, algorithmic characterizations (over simple spaces such as Cantor space and Euclidean space). This was originally shown for Hausdorff dimension in 2017 by J.~Lutz and N.~Lutz \cite{LL18}. 

\begin{theorem}[Point-to-Set Principle; Theorem 1 of \cite{LL18}]
    \label{theorem:PTS_for_H_dimension}
    For any $X \subseteq 2^\omega$ or $X \subseteq \mathbb{R}^m$, we have
    \begin{align*}
        \dim_{\operatorname{H}} X = \min_{B \in 2^{\leq \omega}} \sup_{x \in X} \dim^B(x).
    \end{align*}
\end{theorem}
Notably, this implies that for each subset $X$, there exists an oracle $B \in 2^{\leq \omega}$ which is powerful enough to simulate Hausdorff dimension by ML-$B$-$s$-tests. Any such oracle is called a \textit{Hausdorff oracle} for $X$. An analogous formula holds for classical packing dimension (see Theorem 2 of \cite{LL18}), swapping weak effective dimension $\dim$ for strong dimension: $\operatorname{Dim}$. 

The main insight used to prove the Point-to-Set Principle involves the ability to build for a fixed set $X \subseteq \mathbb{R}^m$ an oracle $B$ with knowledge of how $X$ intersects with any given dyadic cube from $\mathcal{Q}^m$. Arbitrary covers of $X$ exhibiting $\dim_{\operatorname{H}} X$ may be approximated using covers comprising dyadic cubes, and so the Hausdorff dimension of $X$ is the supremum of the size of the intersection with $X$ along nested sequences of dyadic cubes. Suppose that the intersection of a nested sequence intersecting with $X$ is the singleton set containing $x \in X$. Then $\dim^B(x)$ exactly captures how much dimension $X$ takes on at $x$.

% Provide insight into why it's true... Ryan shared his geometric understanding of it too

In practice, the power of the Point-to-Set Principle is in proving lower-bounds for Hausdorff dimension. As remarked to me by Ryan~Bushling, one might as well consider this a \textit{local-to-global} principle for the following reason: in order to witness a lower bound on $\dim_{\operatorname{H}}X$, it is not enough to find a single sufficiently-random element $x \in X$, for any countable number of reals could be coded into an oracle. Instead, one must demonstrate the existence of \textit{many} highly-complex elements of $X$, which requires knowledge of a large portion of the set to do. 

The Point-to-Set Principle has found applications across geometric measure theory, including in bounding the Hausdorff dimension of Furstenberg sets \cite{LS16, Stu18}, lineal extensions \cite{BF25}, distance sets \cite{ABW23}, and pinned-distance sets \cite{Stu22_pinned}; as well as in extending the Marstrand-Mattila Projection Theorem \cite{Stu21}. 

 % Correspondence Principle for original PTS
Well before the full Point-to-Set Principle was proved, L.~Staiger \cite{Sta89} and B.~Ryabko \cite{Rya93, Rya94} had already observed versions of the Point-to-Set Principle for simply-definable classes such as $\Sigma^0_2$-classes. This was slightly generalized by J.~Hitchcock in \cite{Hit05} in what he calls a \textit{correspondence principle} between constructive dimension and Hausdorff dimension \cite{Hit05}. The spirit of any correspondence principle is to characterize a classical notion by unrelativized, effective means (i.e., avoiding oracles). 
 
\begin{theorem}[Correspondence Principle for Hausdorff dimension, Theorem 5.3 of \cite{Hit05}]
    \label{theorem:correspondence_principle_for_H_dimension}
     For any $\Sigma^0_2$-class $X \subseteq 2^\omega$, we have
     \begin{align*}
         \dim_{\operatorname{H}} X = \dim(X).
     \end{align*}
\end{theorem}

Hitchcock further showed that this correspondence broke down immediately beyond this level of the lightface Borel hierarchy. 

\subsection{PTS for Hausdorff Measures}

% Lutz & Miller's PTS principles
The original Point-to-Set Principle has recently been refined to the $s$-dimensional Hausdorff outer measures by P.~Lutz and J.~Miller \cite{LM24}. These results also expand on the complexity characterizations of various partial randomness notions from Theorem~\ref{theorem:complexity_characterizations_of_randomness}.

\begin{theorem} 
    \label{theorem:PTS_log_s_measure_and_sigma_finite}
    For every $X \subseteq 2^\omega$ and $s \geq 0$,
    \begin{align*}
        \log \mathcal{H}^s(X) =^+ \inf_{B \in 2^{\leq \omega}} \sup_{x \in X} \liminf_{r \to \infty}  \operatorname{\Delta M}^B_s(\restrict{x}{r}),
    \end{align*}
    and X is not $\sigma$-finite with respect to $\mathcal{H}^s$ if and only if
    \begin{align*}
        (\forall B \in 2^{\leq \omega})(\exists x \in X) \brk{ \liminf_{r \to \infty} \operatorname{\Delta M}^B_s(\restrict{x}{r}) = \infty }.
    \end{align*}
\end{theorem}

By Theorem~\ref{theorem:complexity_characterizations_of_randomness}, $X$ is not $\sigma$-finite in dimension $s$ if and only if there is a strong Solovay-$B$-$s$-random element in $X$ for every oracle $B$. 

\begin{corollary}
    \label{corollary:PTS_null_or_finite}
    For every $X \subseteq 2^\omega$ and $s \geq 0$,
    \begin{align*}
        \mathcal{H}^s(X) > 0 & \iff(\forall B \in 2^{\leq \omega})(\exists x \in X)\brk{\liminf_{r \to \infty} \operatorname{\Delta M}_s^B(\restrict{x}{r}) > -\infty}\\
        & \iff(\forall B \in 2^{\leq \omega})(\exists x \in X)\brk{\liminf_{r \to \infty} \operatorname{\Delta P}_s^B(\restrict{x}{r}) > -\infty}, \\
        \mathcal{H}^s(X) < \infty &\iff (\exists B \in 2^{\leq \omega}) \brk{\sup_{x \in X} \liminf_{r \to \infty} \operatorname{\Delta M}_s^B(\restrict{x}{r}) < \infty}.
    \end{align*} 
\end{corollary}

Again by Theorem~\ref{theorem:complexity_characterizations_of_randomness}, $X$ is not $s$-null if and only if there is an ML-$B$-$s$-random element in $X$ for every oracle $B$. 

In effect, these refined point-to-set principles demonstrate the usefulness of \textit{a priori} complexity--defined by continuous semimeasures--in characterizing not just Hausdorff dimension but also the $s$-dimensional Hausdorff outer measures. 